\documentclass{article}
\usepackage[margin=0.8in]{geometry}
\usepackage{amsmath,amssymb,amsthm,mathtools}
\usepackage{mathrsfs,bm,bbm}
\usepackage{graphicx,booktabs,array,multirow,tabularx,longtable}
\usepackage{enumitem}
\usepackage{xcolor}
\usepackage{algorithm,algpseudocode}
\usepackage{subcaption,tikz}
\usepackage{siunitx,cancel,accents}
\usepackage{microtype}
\usepackage[numbers,sort&compress]{natbib}
\usepackage[colorlinks=true,linkcolor=blue,citecolor=blue,urlcolor=blue]{hyperref}
\usepackage{cleveref}
\setlength{\emergencystretch}{2em}
\newtheorem{assumption}{Assumption}
\newtheorem{definition}{Definition}
\newtheorem{theorem}{Theorem}
\newtheorem{lemma}{Lemma}
\newtheorem{proposition}{Proposition}
\newtheorem{openendedquestion}{Open-ended Question}
\newtheorem{conjecture}{Conjecture}
\newcommand{\coreref}[1]{\ref{#1}}

\begin{document}

\sloppy

\section*{panel\_\allowbreak{}rarelink\_\allowbreak{}cancellation}

\noindent\textit{Specialization:} v1\par

\paragraph{TL;DR.} For every probability triple in the open unit cube, the normal-link and multiplicative distributional-DiD maps agree exactly on the two planes \(p_1=p_3\) or \(p_2=p_3\), and their complete off-plane discrepancy has a cancellation-safe signed-rectangle factorization. On the comparable-rarity box, the resolved common-cutoff theorem sharpens its coefficient uniformly to \(C_{\rho,\ell}=-u^{-2}+O(u^{-4})\), including both planes. The resulting Gaussian drift frontiers, bounded-count product-Poisson limits, and finitely terminating rational confidence intervals are proved for the deliberately independent sixty-role split design.

\section{Project justification}

\paragraph{Gap.} Djuazon and Tsyawo's Assumptions 1--2 identify the fixed-threshold standard-normal DTT functional, but neither their interior inference theory nor one-dimensional inverse-Mills bounds characterize when that normal-link map equals its multiplicative counterpart over the full probability cube. At the rare boundary, the remaining analytic issue is uniform control of a mixed derivative without division by either cancellation coordinate.

\paragraph{Niche.} The project separates a global population geometry, valid for arbitrary \(p_1,p_2,p_3\in(0,1)\), from a comparable-rarity refinement that links cancellation distance to the Gaussian and bounded-count sampling frontiers.

\paragraph{Fill.} The paper proves the global signed-rectangle factorization and exact two-plane zero set, resolves the uniform common-cutoff expansion on the fixed coefficient box, and propagates that refinement through pairwise and aggregate Gaussian limits, a product-Poisson transition, and certified finite-sample outer inference.

\section{Related work}

Djuazon and Tsyawo (2026) supply, through their Assumptions 1--2 and Theorem 1, the causal identification baseline specialized here to a standard-normal link and threshold zero; that population identity is not a novelty claim. Their interior Theorem 2 motivates comparison with regular Gaussian inference, but the displayed split-role proposition is proved independently for sixty mutually independent role samples and is an analogue rather than a claim of membership in their original sampling class. Pinelis (2019) and Baricz (2008) develop one-dimensional inverse-Mills tools, whereas the global signed-rectangle identity and its exact two-plane cancellation set are the present paper's mixed construction. Chen and Fang (2019) explain higher-order behavior under degeneracy in Gaussian experiments, and Marohn (2001) records the distinct Poisson limit generated by bounded rare counts. Sachs, Gabriel, and Fay (2026) motivate exact inversion for nonlinear functions of discrete probabilities; here the rational outer algorithm is proved to terminate and to retain simultaneous coverage, without an optimality comparison. Fernández-Val, Meier, van Vuuren, and Vella (2024), Wooldridge (2023), and Roth and Sant'Anna (2023) emphasize that nonlinear link choice is substantive in difference-in-differences models.

\section{Setup}

Target (estimand / structural parameter): \(\overline\Delta=15^{-1}\sum_{\ell=1}^{15}\{\Pr(Y_{t_{\ell}}(1)\le0\mid G=g_{\ell})-\Pr(Y_{t_{\ell}}(0)\le0\mid G=g_{\ell})\}\).
Primary functional / estimator / diagnostic: Under normal-link distributional parallel trends and no anticipation, the causal estimand is identified by the observed-law functional \(\overline\Delta=15^{-1}\sum_{\ell=1}^{15}\{H(r_{\ell1},r_{\ell2},r_{\ell3})-r_{\ell4}\}\). On the common-rarity class this equals \(15^{-1}\sum_{\ell=1}^{15}\{\Phi[\Phi^{-1}(\rho c_{\ell1})+\Phi^{-1}(\rho c_{\ell2})-\Phi^{-1}(\rho c_{\ell3})]-\rho c_{\ell4}\}\). Its zero-safe observed-count plug-in is \(15^{-1}\sum_{\ell=1}^{15}\widehat\Delta^E_{\ell}=15^{-1}\sum_{\ell=1}^{15}\{H(\widetilde r_{\ell1},\widetilde r_{\ell2},\widetilde r_{\ell3})-\widetilde r_{\ell4}\}\), where \(\widetilde r_{\ell j}=(K_{\ell j}+1/2)/(n_{\ell j}+1)\)..
\begin{itemize}
  \item \texttt{L} --- \texttt{positive integer}; \texttt{\textbackslash{}\allowbreak{}mathbb N}; \texttt{fixed dimension}
  \par\smallskip\noindent\emph{Definition.} \(L=15\) is the fixed number of pre/post comparisons.
  \item \texttt{J} --- \texttt{finite index set}; \(\{1,2,3,4\}\); \texttt{role index}
  \par\smallskip\noindent\emph{Definition.} \(J=\{1,2,3,4\}\) indexes treated-pre, control-post, control-pre, and treated-post roles.
  \item \texttt{alpha} --- \texttt{scalar}; \((0,1)\); \texttt{inference constant}
  \par\smallskip\noindent\emph{Definition.} \(\alpha\) is the simultaneous noncoverage level.
  \item \texttt{beta} --- \texttt{scalar}; \((0,1)\); \texttt{inference constant}
  \par\smallskip\noindent\emph{Definition.} \(\beta\) is a pairwise noncoverage budget, set to \(\alpha/L\) for aggregation.
  \item \texttt{eta} --- \texttt{positive rational scalar}; \(\mathbb Q_{>0}\); \texttt{algorithmic tolerance}
  \par\smallskip\noindent\emph{Definition.} \(\eta\) is the requested rational numerical enclosure tolerance.
  \item \texttt{n} --- \texttt{positive integer sequence}; \(\mathbb N\); \texttt{sample-\allowbreak{}size scale}
  \par\smallskip\noindent\emph{Definition.} \(n\) is the common asymptotic scale for all role samples.
  \item \texttt{nrole} --- \texttt{array of positive integers}; \(\mathbb N^{L\times4}\); \texttt{role sample sizes}
  \par\smallskip\noindent\emph{Definition.} \(n_{\ell j}\) is the number of predeclared blocks assigned to pair \(\ell\) and role \(j\).
  \item \texttt{pi\_\allowbreak{}lower} --- \texttt{positive scalar}; \((0,\infty)\); \texttt{sampling constant}
  \par\smallskip\noindent\emph{Definition.} \(\underline\pi\) is the fixed lower role-share bound.
  \item \texttt{pi\_\allowbreak{}upper} --- \texttt{positive scalar}; \((0,\infty)\); \texttt{sampling constant}
  \par\smallskip\noindent\emph{Definition.} \(\overline\pi\) is the fixed upper role-share bound.
  \item \texttt{pi} --- \texttt{positive array}; \((0,\infty)^{L\times4}\); \texttt{asymptotic sampling shares}
  \par\smallskip\noindent\emph{Definition.} \(\pi_{\ell j}\) is the candidate limiting share for pair \(\ell\) and role \(j\).
  \item \texttt{Pairs} --- \texttt{fixed ordered collection}; \texttt{panel comparison design}
  \par\smallskip\noindent\emph{Definition.} \(\mathcal P=\{(g_{\ell},h_{\ell},s_{\ell},t_{\ell}):1\le\ell\le L\}\) is the predeclared list of treated groups, comparison groups, pre-periods, and post-periods.
  \item \texttt{Y} --- \texttt{potential-\allowbreak{}outcome array}; \texttt{causal primitive}
  \par\smallskip\noindent\emph{Definition.} \(Y_{i\nu}(d)\) is unit \(i\)'s outcome at time \(\nu\) under treatment state \(d\in\{0,1\}\).
  \item \texttt{Ggrp} --- \texttt{group label}; \texttt{causal primitive}
  \par\smallskip\noindent\emph{Definition.} \(G_i\) is unit \(i\)'s treatment-cohort or comparison-group label.
  \item \texttt{y0} --- \texttt{outcome threshold}; \texttt{target index}
  \par\smallskip\noindent\emph{Definition.} \(y_0=0\) is the fixed crime-free threshold.
  \item \texttt{Phi} --- \texttt{function}; \(\mathbb R\to(0,1)\); \texttt{working link}
  \par\smallskip\noindent\emph{Definition.} \(\Phi\) is the standard-normal distribution function.
  \item \texttt{phi} --- \texttt{function}; \(\mathbb R\to\mathbb R_{>0}\); \texttt{tail derivative}
  \par\smallskip\noindent\emph{Definition.} \(\phi\) is the standard-normal density.
  \item \texttt{q} --- \texttt{function}; \((0,1)\to\mathbb R\); \texttt{inverse link}
  \par\smallskip\noindent\emph{Definition.} \(q\) is the standard-normal quantile function.
  \item \texttt{rho} --- \texttt{positive scalar sequence}; \((0,1)\); \texttt{structural regime parameter}
  \par\smallskip\noindent\emph{Definition.} \(\rho=\rho_n\) is the common exceedance-rarity scale.
  \item \texttt{c} --- \texttt{positive coefficient array}; \((0,\infty)^{L\times4}\); \texttt{primitive cell coefficients}
  \par\smallskip\noindent\emph{Definition.} \(c=(c_{\ell j,n})\) contains the relative exceedance coefficients.
  \item \texttt{r} --- \texttt{probability array}; \((0,1)^{L\times4}\); \texttt{observed-\allowbreak{}data primitive}
  \par\smallskip\noindent\emph{Definition.} For \(\mathcal P_\ell=(g_\ell,h_\ell,s_\ell,t_\ell)\), \(r_{\ell1}=\Pr(Y_{s_\ell}(1)>y_0\mid G=g_\ell)\), \(r_{\ell2}=\Pr(Y_{t_\ell}(0)>y_0\mid G=h_\ell)\), \(r_{\ell3}=\Pr(Y_{s_\ell}(0)>y_0\mid G=h_\ell)\), and \(r_{\ell4}=\Pr(Y_{t_\ell}(1)>y_0\mid G=g_\ell)\).
  \item \texttt{Erole} --- \texttt{binary triangular array}; \(\{0,1\}^{L\times4\times\mathbb N}\); \texttt{role-\allowbreak{}level observation}
  \par\smallskip\noindent\emph{Definition.} \(E_{\ell j a,n}\) is the threshold-exceedance indicator for observation \(a\) in the predeclared sample for pair \(\ell\) and role \(j\).
  \item \texttt{K} --- \texttt{count array}; \(\prod_{\ell,j}\{0,\ldots,n_{\ell j}\}\); \texttt{observed statistic}
  \par\smallskip\noindent\emph{Definition.} \(K_{\ell j,n}=\sum_{a=1}^{n_{\ell j}}E_{\ell j a,n}\) counts threshold exceedances in pair \(\ell\), role \(j\).
  \item \texttt{H} --- \texttt{function}; \((0,1)^3\to(0,1)\); \texttt{identified counterfactual}
  \par\smallskip\noindent\emph{Definition.} \(H\) is the exact normal-link counterfactual exceedance map.
  \item \texttt{Pmap} --- \texttt{function}; \((0,1)^3\to\mathbb R_{>0}\); \texttt{rare-\allowbreak{}event comparator}
  \par\smallskip\noindent\emph{Definition.} \(P^{\rm P}\) is the Poisson-ratio counterfactual map.
  \item \texttt{Q} --- \texttt{positive array}; \texttt{Poisson coefficient}
  \par\smallskip\noindent\emph{Definition.} \(Q_{\ell}=c_{\ell1}c_{\ell2}/c_{\ell3}\).
  \item \texttt{x} --- \texttt{real array}; \texttt{first cancellation coordinate}
  \par\smallskip\noindent\emph{Definition.} \(x_{\ell}=\log(c_{\ell1}/c_{\ell3})\).
  \item \texttt{z} --- \texttt{real array}; \texttt{second cancellation coordinate}
  \par\smallskip\noindent\emph{Definition.} \(z_{\ell}=\log(c_{\ell2}/c_{\ell3})\).
  \item \texttt{A} --- \texttt{real array}; \texttt{cancellation product}
  \par\smallskip\noindent\emph{Definition.} \(A_{\ell}=x_{\ell}z_{\ell}\).
  \item \texttt{u} --- \texttt{positive scalar}; \texttt{common inverse-\allowbreak{}tail cutoff}
  \par\smallskip\noindent\emph{Definition.} \(u=-q(\rho)\).
  \item \texttt{m} --- \texttt{function}; \(\mathbb R\to\mathbb R_{>0}\); \texttt{analytic coefficient}
  \par\smallskip\noindent\emph{Definition.} \(m\) is the left-tail inverse Mills ratio.
  \item \texttt{B} --- \texttt{function with rare-\allowbreak{}panel evaluations}; \((0,1)^3\to\mathbb R\); \texttt{mixed log-\allowbreak{}error coefficient}
  \par\smallskip\noindent\emph{Definition.} \(B:(0,1)^3\to\mathbb R\) is the mixed log-error coefficient. For rare-panel arguments, \(B_{\rho,\ell}=B(\rho c_{\ell1},\rho c_{\ell2},\rho c_{\ell3})\).
  \item \texttt{psi} --- \texttt{function}; \(\mathbb R\to\mathbb R_{>0}\); \texttt{cancellation-\allowbreak{}safe exponential quotient}
  \par\smallskip\noindent\emph{Definition.} \(\psi(0)=1\) and \(\psi(w)=(e^w-1)/w\) for \(w\ne0\).
  \item \texttt{Ccoef} --- \texttt{function with rare-\allowbreak{}panel evaluations}; \((0,1)^3\to\mathbb R\); \texttt{complete-\allowbreak{}error coefficient}
  \par\smallskip\noindent\emph{Definition.} For \(p_1,p_2,p_3\in(0,1)\), \(C(p_1,p_2,p_3)=B(p_1,p_2,p_3)\psi\!\left(\left[\log(p_1/p_3)\log(p_2/p_3)\right]B(p_1,p_2,p_3)\right)\). For rare-panel arguments, \(C_{\rho,\ell}=C(\rho c_{\ell1},\rho c_{\ell2},\rho c_{\ell3})\).
  \item \texttt{Delta} --- \texttt{real array}; \texttt{pairwise causal DTT}
  \par\smallskip\noindent\emph{Definition.} \(\Delta_{\ell}=\Pr(Y_{t_{\ell}}(1)\le y_0\mid G=g_{\ell})-\Pr(Y_{t_{\ell}}(0)\le y_0\mid G=g_{\ell})\) is the pairwise crime-free DTT.
  \item \texttt{DeltaP} --- \texttt{real array}; \texttt{Poisson approximation target}
  \par\smallskip\noindent\emph{Definition.} \(\Delta^P_{\ell}=P^{\rm P}(r_{\ell1},r_{\ell2},r_{\ell3})-r_{\ell4}=\rho(Q_{\ell}-c_{\ell4})\).
  \item \texttt{barDelta} --- \texttt{real scalar}; \texttt{aggregate causal estimand}
  \par\smallskip\noindent\emph{Definition.} \(\overline\Delta=L^{-1}\sum_{\ell=1}^{L}\Delta_{\ell}\).
  \item \texttt{rtilde} --- \texttt{probability array}; \texttt{zero-\allowbreak{}safe cell estimator}
  \par\smallskip\noindent\emph{Definition.} \(\widetilde r_{\ell j}=(K_{\ell j}+1/2)/(n_{\ell j}+1)\).
  \item \texttt{hatDeltaE} --- \texttt{real array}; \texttt{exact-\allowbreak{}link estimator}
  \par\smallskip\noindent\emph{Definition.} \(\widehat\Delta^E_{\ell}=H(\widetilde r_{\ell1},\widetilde r_{\ell2},\widetilde r_{\ell3})-\widetilde r_{\ell4}\).
  \item \texttt{hatDeltaP} --- \texttt{real array}; \texttt{Poisson-\allowbreak{}ratio estimator}
  \par\smallskip\noindent\emph{Definition.} \(\widehat\Delta^P_{\ell}=P^{\rm P}(\widetilde r_{\ell1},\widetilde r_{\ell2},\widetilde r_{\ell3})-\widetilde r_{\ell4}\).
  \item \texttt{V} --- \texttt{positive array}; \texttt{pairwise rare Gaussian variance}
  \par\smallskip\noindent\emph{Definition.} \(V_{\ell}=c_{\ell4}/\pi_{\ell4}+Q_{\ell}^{2}\sum_{j=1}^{3}(\pi_{\ell j}c_{\ell j})^{-1}\).
  \item \texttt{barV} --- \texttt{positive scalar}; \texttt{aggregate rare Gaussian variance}
  \par\smallskip\noindent\emph{Definition.} \(\overline V=L^{-2}\sum_{\ell=1}^{L}V_{\ell}\).
  \item \texttt{tau} --- \texttt{nonnegative scalar sequence}; \texttt{experiment frontier parameter}
  \par\smallskip\noindent\emph{Definition.} \(\tau_n=n\rho_n\) is the expected-count scale.
  \item \texttt{Lambda} --- \texttt{real array}; \texttt{pairwise drifting-\allowbreak{}cancellation index}
  \par\smallskip\noindent\emph{Definition.} \(\Lambda_{\ell n}=A_{\ell n}\sqrt{\tau_n}/u_n^2\).
  \item \texttt{bshift} --- \texttt{real scalar sequence}; \texttt{exact aggregate approximation shift}
  \par\smallskip\noindent\emph{Definition.} \(b_n=\sqrt{n/\rho_n}\{L\sqrt{\overline V}\}^{-1}\sum_{\ell=1}^{L}(\Delta^P_{\ell n}-\Delta_{\ell n})\).
  \item \texttt{Z} --- \texttt{Poisson array}; \texttt{limit experiment}
  \par\smallskip\noindent\emph{Definition.} \(Z_{\ell j}\sim\operatorname{Poisson}(\pi_{\ell j}\tau c_{\ell j})\) independently in the bounded-count limit.
  \item \texttt{W} --- \texttt{positive array}; \texttt{smoothed Poisson cell limit}
  \par\smallskip\noindent\emph{Definition.} \(W_{\ell j}=(Z_{\ell j}+1/2)/\pi_{\ell j}\).
  \item \texttt{T} --- \texttt{real array}; \texttt{product-\allowbreak{}Poisson ratio limit}
  \par\smallskip\noindent\emph{Definition.} \(T_{\ell}=W_{\ell1}W_{\ell2}/W_{\ell3}-W_{\ell4}\).
  \item \texttt{Pcone} --- \texttt{compact subset}; \([0,1/4]^4\); \texttt{finite-\allowbreak{}sample comparable-\allowbreak{}rarity closure}
  \par\smallskip\noindent\emph{Definition.} \(\mathcal P_0=\{p\in[0,1/4]^4:p_i\le8p_j\text{ for every }i,j\}\).
  \item \texttt{Gmap} --- \texttt{function}; \(\mathcal P_0\to\mathbb R\); \texttt{confidence target}
  \par\smallskip\noindent\emph{Definition.} \(G\) is the rare-link DTT map on the closed cone.
  \item \texttt{Llip} --- \texttt{positive integer}; \texttt{global max-\allowbreak{}norm Lipschitz certificate}
  \par\smallskip\noindent\emph{Definition.} \(L_{\rm lip}=2^{258}\).
  \item \texttt{Ddev} --- \texttt{extended-\allowbreak{}real function}; \texttt{finite-\allowbreak{}sample deviance}
  \par\smallskip\noindent\emph{Definition.} \(D_p(k)=2\sum_{j=1}^{4}n_j\operatorname{KL}(\operatorname{Bern}(k_j/n_j),\operatorname{Bern}(p_j))\), with endpoint conventions.
  \item \texttt{Rregion} --- \texttt{set-\allowbreak{}valued function}; \texttt{outer probability region}
  \par\smallskip\noindent\emph{Definition.} \(\mathcal R_{\beta}(k)=\{p\in\mathcal P_0:D_p(k)\le8/\beta\}\).
  \item \texttt{Alg} --- \texttt{deterministic algorithm}; \texttt{certified interval map}
  \par\smallskip\noindent\emph{Definition.} \(\mathsf{Outer}_{\beta,\eta}\) is the finite rational-grid construction defined below.
  \item \texttt{Iinterval} --- \texttt{rational closed interval}; \texttt{pairwise confidence interval}
  \par\smallskip\noindent\emph{Definition.} \(I_{\ell}(K)=\mathsf{Outer}_{\beta,\eta}(K_{\ell1},\ldots,K_{\ell4})\).
  \item \texttt{barI} --- \texttt{rational closed interval}; \texttt{aggregate confidence interval}
  \par\smallskip\noindent\emph{Definition.} If \(I_{\ell}=[L_{\ell},U_{\ell}]\), then \(\overline I=[L^{-1}\sum_{\ell}L_{\ell},L^{-1}\sum_{\ell}U_{\ell}]\).
  \item \texttt{epsilon\_\allowbreak{}int} --- \texttt{positive scalar}; \((0,1/2)\); \texttt{interior benchmark constant}
  \par\smallskip\noindent\emph{Definition.} \(\varepsilon_{\rm int}\) is the fixed interior-probability margin used only for the published-theorem reduction.
  \item \texttt{pint} --- \texttt{probability array}; \((0,1)^{L\times4}\); \texttt{interior benchmark probability}
  \par\smallskip\noindent\emph{Definition.} \(p^{\circ}_{\ell j}\) is the fixed role probability in the interior benchmark experiment.
  \item \texttt{gamma\_\allowbreak{}int} --- \texttt{real array}; \texttt{interior exact-\allowbreak{}link gradient}
  \par\smallskip\noindent\emph{Definition.} Writing \(v^{\circ}_{\ell}=q(p^{\circ}_{\ell1})+q(p^{\circ}_{\ell2})-q(p^{\circ}_{\ell3})\), define \(\gamma^{\circ}_{\ell1}=\phi(v^{\circ}_{\ell})/\phi(q(p^{\circ}_{\ell1}))\), \(\gamma^{\circ}_{\ell2}=\phi(v^{\circ}_{\ell})/\phi(q(p^{\circ}_{\ell2}))\), \(\gamma^{\circ}_{\ell3}=-\phi(v^{\circ}_{\ell})/\phi(q(p^{\circ}_{\ell3}))\), and \(\gamma^{\circ}_{\ell4}=-1\).
  \item \texttt{IFint} --- \texttt{role-\allowbreak{}indexed function array}; \texttt{interior split-\allowbreak{}role influence function}
  \par\smallskip\noindent\emph{Definition.} \(\iota^{\circ}_{\ell j}(e)=\gamma^{\circ}_{\ell j}(e-p^{\circ}_{\ell j})/\pi_{\ell j}\) for \(e\in\{0,1\}\).
  \item \texttt{Omega\_\allowbreak{}int} --- \texttt{positive array}; \texttt{interior split-\allowbreak{}role asymptotic variance}
  \par\smallskip\noindent\emph{Definition.} \(\Omega^{\circ}_{\ell}=\sum_{j=1}^{4}(\gamma^{\circ}_{\ell j})^2p^{\circ}_{\ell j}(1-p^{\circ}_{\ell j})/\pi_{\ell j}\).
\end{itemize}

\section{Assumptions}

\begin{assumption}[ass:normal-link-parallel-trends]\label{ass:normal-link-parallel-trends}
For every \(\ell\), \(q(\Pr(Y_{t_{\ell}}(0)>y_0\mid G=g_{\ell}))-q(\Pr(Y_{s_{\ell}}(0)>y_0\mid G=g_{\ell}))=q(\Pr(Y_{t_{\ell}}(0)>y_0\mid G=h_{\ell}))-q(\Pr(Y_{s_{\ell}}(0)>y_0\mid G=h_{\ell}))\).
\par\smallskip\noindent\emph{Standard} (Normal-link distributional parallel trends, \cite{DjuazonTsyawo2026}).
\end{assumption}

\begin{assumption}[ass:no-anticipation]\label{ass:no-anticipation}
For every \(\ell\), \(\Pr(Y_{s_{\ell}}(1)>y_0\mid G=g_{\ell})=\Pr(Y_{s_{\ell}}(0)>y_0\mid G=g_{\ell})\).
\par\smallskip\noindent\emph{Standard} (Distributional no anticipation, \cite{DjuazonTsyawo2026}).
\end{assumption}

\begin{assumption}[ass:coefficient-box]\label{ass:coefficient-box}
For every \(\ell,j,n\), \(c_{\ell j,n}\in[1/2,4]\).
\par\smallskip\noindent\emph{Novel.} Comparable rarity allows persistent group and period heterogeneity by a factor of eight and is realistic for threshold exceedances across a fixed collection of nearby blocks and months.
\end{assumption}

\begin{assumption}[ass:rarity-scale]\label{ass:rarity-scale}
For every \(n\), \(0<\rho_n\le1/16\).
\par\smallskip\noindent\emph{Novel.} The condition describes a sparse-event sequence while keeping every cell probability in the proper normal-link domain; crime incidence and other fixed-threshold count events provide concrete examples.
\end{assumption}

\begin{assumption}[ass:common-rarity]\label{ass:common-rarity}
For every \(\ell,j,n\), \(r_{\ell j,n}=\rho_n c_{\ell j,n}\).
\par\smallskip\noindent\emph{Novel.} A common exposure intensity with bounded cell-specific multipliers is natural when rare exceedances share a citywide or periodwide baseline hazard.
\end{assumption}

\begin{assumption}[ass:bernoulli-role-sampling]\label{ass:bernoulli-role-sampling}
For every \(\ell,j,n\) and \(1\le a\le n_{\ell j}\), \(E_{\ell j a,n}\sim\operatorname{Bernoulli}(r_{\ell j,n})\).
\par\smallskip\noindent\emph{Standard} (Bernoulli sampling of threshold indicators, \cite{BrownCaiDasGupta2001}).
\end{assumption}

\begin{assumption}[ass:disjoint-role-independence]\label{ass:disjoint-role-independence}
The collection \(\{E_{\ell j a,n}:1\le\ell\le L, j\in J,1\le a\le n_{\ell j}\}\) is mutually independent for every \(n\).
\par\smallskip\noindent\emph{Novel.} Outcome-independent splitting of independent complete blocks into predeclared pair--role samples realizes this condition while preserving each role's marginal event probability.
\end{assumption}

\begin{assumption}[ass:role-share-bounds]\label{ass:role-share-bounds}
For every \(\ell,j,n\), \(\underline\pi\le n_{\ell j}/n\le\overline\pi\).
\par\smallskip\noindent\emph{Standard} (Positive sampling proportions, \cite{DjuazonTsyawo2026}).
\end{assumption}

\begin{assumption}[ass:role-share-limits]\label{ass:role-share-limits}
For every \(\ell,j\), \(n_{\ell j}/n\to\pi_{\ell j}\).
\par\smallskip\noindent\emph{Standard} (Convergent sampling proportions, \cite{DjuazonTsyawo2026}).
\end{assumption}

\begin{assumption}[ass:coefficient-limits]\label{ass:coefficient-limits}
For every \(\ell,j\), \(c_{\ell j,n}\to c_{\ell j}\).
\par\smallskip\noindent\emph{Novel.} Stable relative rare-event risks across a fixed set of cells are compatible with triangular arrays generated by a shrinking common exposure intensity.
\end{assumption}

\begin{assumption}[ass:rare-limit]\label{ass:rare-limit}
\(\rho_n\to0\).
\par\smallskip\noindent\emph{Novel.} The sequence formalizes increasingly rare fixed-threshold events and is the boundary excluded by interior normal-link inference.
\end{assumption}

\begin{assumption}[ass:interior-fixed-probabilities]\label{ass:interior-fixed-probabilities}
For every \(\ell,j,n\), \(r_{\ell j,n}=p^{\circ}_{\ell j}\).
\par\smallskip\noindent\emph{Standard} (Fixed interior probability experiment, \cite{DjuazonTsyawo2026}).
\end{assumption}

\begin{assumption}[ass:interior-probability-margin]\label{ass:interior-probability-margin}
For every \(\ell,j\), \(\varepsilon_{\rm int}\le p^{\circ}_{\ell j}\le1-\varepsilon_{\rm int}\).
\par\smallskip\noindent\emph{Standard} (Djuazon--Tsyawo Assumption 5 specialized to a finite threshold and standard-normal link, \cite{DjuazonTsyawo2026}).
\end{assumption}

\section{Definitions}

\begin{definition}[def:rare-panel-class]\label{def:rare-panel-class}
\par\smallskip\noindent\emph{Defined object.} \texttt{Comparable-\allowbreak{}rarity split-\allowbreak{}panel class}
\par\smallskip\noindent\emph{Construction.} \(\mathcal M_n=\{(Y,G,\mathcal P,\rho_n,c_n,n_{\ell j},E_n,K_n):\text{the member satisfies assumptions }\mathrm{ass:normal\mbox{-}link\mbox{-}parallel\mbox{-}trends},\mathrm{ass:no\mbox{-}anticipation},\mathrm{ass:coefficient\mbox{-}box},\mathrm{ass:rarity\mbox{-}scale},\mathrm{ass:common\mbox{-}rarity},\mathrm{ass:bernoulli\mbox{-}role\mbox{-}sampling},\mathrm{ass:disjoint\mbox{-}role\mbox{-}independence},\mathrm{ass:role\mbox{-}share\mbox{-}bounds}\}\).
\par\smallskip\noindent\emph{Member properties.} \texttt{ass:\allowbreak{}normal-\allowbreak{}link-\allowbreak{}parallel-\allowbreak{}trends}; \texttt{ass:\allowbreak{}no-\allowbreak{}anticipation}; \texttt{ass:\allowbreak{}coefficient-\allowbreak{}box}; \texttt{ass:\allowbreak{}rarity-\allowbreak{}scale}; \texttt{ass:\allowbreak{}common-\allowbreak{}rarity}; \texttt{ass:\allowbreak{}bernoulli-\allowbreak{}role-\allowbreak{}sampling}; \texttt{ass:\allowbreak{}disjoint-\allowbreak{}role-\allowbreak{}independence}; \texttt{ass:\allowbreak{}role-\allowbreak{}share-\allowbreak{}bounds}.
\end{definition}

\begin{definition}[def:mixed-coefficient]\label{def:mixed-coefficient}
\par\smallskip\noindent\emph{Defined object.} \texttt{Mixed log-\allowbreak{}error coefficient}
\par\smallskip\noindent\emph{Construction.} For any \(p=(p_1,p_2,p_3)\in(0,1)^3\), put \(b=p_3\), \(x=\log(p_1/p_3)\), \(z=\log(p_2/p_3)\), and \(A=xz\), and define
\[
B(p_1,p_2,p_3)=\int_0^1\!\int_0^1
\frac{m'\!\left(q(be^{sx})+q(be^{tz})-q(b)\right)}
{m\!\left(q(be^{sx})\right)m\!\left(q(be^{tz})\right)}\,dt\,ds,
\qquad
C(p_1,p_2,p_3)=B(p_1,p_2,p_3)\psi\!\left(A B(p_1,p_2,p_3)\right).
\]
For the rare-panel evaluations, write
\[
B_{\rho,\ell}=B(\rho c_{\ell1},\rho c_{\ell2},\rho c_{\ell3}),
\qquad
C_{\rho,\ell}=C(\rho c_{\ell1},\rho c_{\ell2},\rho c_{\ell3}).
\]
\par\smallskip\noindent\emph{Inputs.} \texttt{m}; \texttt{q}; \texttt{psi}; \texttt{rho}; \texttt{c}.
\end{definition}

\begin{definition}[def:frontier-indices]\label{def:frontier-indices}
\par\smallskip\noindent\emph{Defined object.} \texttt{Drifting-\allowbreak{}cancellation indices}
\par\smallskip\noindent\emph{Construction.} The pairwise index is \(\Lambda_{\ell n}=A_{\ell n}\sqrt{n\rho_n}/u_n^2\); the aggregate index is \(b_n=\sqrt{n/\rho_n}\{L\sqrt{\overline V}\}^{-1}\sum_{\ell=1}^{L}(\Delta^P_{\ell n}-\Delta_{\ell n})\).
\par\smallskip\noindent\emph{Inputs.} \texttt{A}; \texttt{n}; \texttt{rho}; \texttt{u}; \texttt{L}; \texttt{barV}; \texttt{DeltaP}; \texttt{Delta}.
\end{definition}

\begin{definition}[def:bounded-count-limit]\label{def:bounded-count-limit}
\par\smallskip\noindent\emph{Defined object.} \texttt{Product-\allowbreak{}Poisson limit statistic}
\par\smallskip\noindent\emph{Construction.} For independent \(Z_{\ell j}\sim\operatorname{Poisson}(\pi_{\ell j}\tau c_{\ell j})\), set \(W_{\ell j}=(Z_{\ell j}+1/2)/\pi_{\ell j}\) and \(T_{\ell}=W_{\ell1}W_{\ell2}/W_{\ell3}-W_{\ell4}\).
\par\smallskip\noindent\emph{Inputs.} \texttt{pi}; \texttt{tau}; \texttt{c}.
\end{definition}

\begin{definition}[def:uniform-proof-handle]\label{def:uniform-proof-handle}
\par\smallskip\noindent\emph{Defined object.} \texttt{Common-\allowbreak{}cutoff inverse-\allowbreak{}tail handle}
\par\smallskip\noindent\emph{Construction.} Start from \(\Phi(-v)/\phi(v)=v^{-1}\int_0^{\infty}e^{-s}e^{-s^2/(2v^2)}\,ds\), retain the signed-rectangle representation for \(B_{\rho,\ell}\), and seek one common cutoff and remainder ledger valid for all \((c_{\ell1},c_{\ell2},c_{\ell3})\in[1/2,4]^3\).
\par\smallskip\noindent\emph{Inputs.} \texttt{Phi}; \texttt{phi}; \texttt{B}; \texttt{Ccoef}; \texttt{c}.
\end{definition}

\begin{definition}[def:rational-algorithm]\label{def:rational-algorithm}
\par\smallskip\noindent\emph{Defined object.} \texttt{Finite rational outer-\allowbreak{}interval algorithm}
\par\smallskip\noindent\emph{Construction.} Given counts \(k\), compute \(D_p(k)\), \(\mathcal R_{\beta}(k)\), and the saturated likelihood \(L_{\rm sat}\); choose a positive rational \(t_{\rm lo}\) with \(L_{\rm sat}e^{-4/\beta}/2\le t_{\rm lo}\le L_{\rm sat}e^{-4/\beta}\), partition \([0,1/4]^4\) into rational boxes of max-diameter \(h\) with \(L_{\rm lip}h\le\eta/8\) and \(\sum_j n_jh\le t_{\rm lo}/4\), retain every box whose rational likelihood upper enclosure meets \(t_{\rm lo}\) and whose intersection with \(\mathcal P_0\) is nonempty, enclose \(G\) at a rational feasible point of each retained box, enlarge by \(L_{\rm lip}h\), and return the rational hull; return \([0,0]\) when the retained list is empty.
\par\smallskip\noindent\emph{Inputs.} \texttt{K}; \texttt{nrole}; \texttt{beta}; \texttt{eta}; \texttt{Llip}; \texttt{Pcone}; \texttt{Gmap}.
\end{definition}

\section{Main results}

\begin{proposition}[prop:published-identification-baseline]\label{prop:published-identification-baseline}
Djuazon and Tsyawo's Theorem 1, specialized to the standard-normal working distribution, threshold \(y_0=0\), and the fixed pairs \(\mathcal P\), identifies \(\Delta_{\ell}=H(r_{\ell1},r_{\ell2},r_{\ell3})-r_{\ell4}\) for every \(\ell\), and hence identifies \(\overline\Delta=L^{-1}\sum_{\ell=1}^{L}\Delta_{\ell}\). Outcome-independent splitting changes only the samples used to estimate the four observed marginals and does not alter this population identity.
\end{proposition}
\begin{proof}
Fix \(\ell\in\{1,\ldots,L\}\), and let
\[
p^0_{\ell t}:=\Pr\!\left(Y_{t_\ell}(0)>y_0\mid G=g_\ell\right).
\tag{1}
\]
The normal-link parallel-trends assumption contains \(q(p^0_{\ell t})\), so its well-typedness, together with the declared domain \((0,1)\) of \(q\), implies \(p^0_{\ell t}\in(0,1)\).  Likewise, the declared support of \(r\) gives \(r_{\ell j}\in(0,1)\) for every \(j\in J\), so every inverse-normal transform below is defined.

By no anticipation and then by the definition of \(r_{\ell1}\),
\[
\Pr\!\left(Y_{s_\ell}(0)>y_0\mid G=g_\ell\right)
=\Pr\!\left(Y_{s_\ell}(1)>y_0\mid G=g_\ell\right)
=r_{\ell1}.
\tag{2}
\]
The definitions of the control-post and control-pre observed marginals give
\[
\Pr\!\left(Y_{t_\ell}(0)>y_0\mid G=h_\ell\right)=r_{\ell2},
\qquad
\Pr\!\left(Y_{s_\ell}(0)>y_0\mid G=h_\ell\right)=r_{\ell3}.
\tag{3}
\]
Substituting (1)--(3) into normal-link parallel trends yields
\[
q(p^0_{\ell t})-q(r_{\ell1})
=q(r_{\ell2})-q(r_{\ell3}),
\]
and hence
\[
q(p^0_{\ell t})
=q(r_{\ell1})+q(r_{\ell2})-q(r_{\ell3}).
\tag{4}
\]
Because \(q=\Phi^{-1}\) on \((0,1)\), applying \(\Phi\) to both sides of (4) gives
\[
p^0_{\ell t}
=\Phi\!\left(q(r_{\ell1})+q(r_{\ell2})-q(r_{\ell3})\right)
=H(r_{\ell1},r_{\ell2},r_{\ell3}).
\tag{5}
\]
Thus (5) identifies the unobserved treated-group untreated-post exceedance probability from three observed marginals; it is a consequence of the two causal assumptions, not a definition of that counterfactual probability.

The definition of \(r_{\ell4}\) gives
\[
\Pr\!\left(Y_{t_\ell}(1)>y_0\mid G=g_\ell\right)=r_{\ell4}.
\tag{6}
\]
For either treatment state, the events \(\{Y_{t_\ell}(d)\le y_0\}\) and \(\{Y_{t_\ell}(d)>y_0\}\) are complements.  Therefore, using the causal definition of \(\Delta_\ell\), followed by (1), (5), and (6),
\[
\begin{aligned}
\Delta_\ell
&=\Pr\!\left(Y_{t_\ell}(1)\le y_0\mid G=g_\ell\right)
  -\Pr\!\left(Y_{t_\ell}(0)\le y_0\mid G=g_\ell\right)\\
&=(1-r_{\ell4})-(1-p^0_{\ell t})\\
&=p^0_{\ell t}-r_{\ell4}\\
&=H(r_{\ell1},r_{\ell2},r_{\ell3})-r_{\ell4}.
\end{aligned}
\tag{7}
\]
Every quantity on the last line of (7) is a functional of the observed marginal law, whereas the first line is the causal target; equations (2)--(7) prove their equality.

Since \(\ell\) was arbitrary, (7) holds for all \(1\le \ell\le L\).  The definition of \(\overline\Delta\) then gives the identifying map
\[
\overline\Delta
=\frac{1}{L}\sum_{\ell=1}^{L}\Delta_\ell
=\frac{1}{L}\sum_{\ell=1}^{L}
\left\{H(r_{\ell1},r_{\ell2},r_{\ell3})-r_{\ell4}\right\}.
\tag{8}
\]
Finally, an outcome-independent split merely assigns observations to role-specific estimation samples without changing any population conditional probability appearing in (1)--(8).  Consequently it can change the estimators and their sampling law, but not the population identities (7)--(8).
\end{proof}
\par\smallskip\noindent \textit{Justification.} This is the fixed-threshold standard-normal implication of Djuazon and Tsyawo's Assumptions 1--2 and fixes the causal estimand before any approximation or inference step. \textit{Gap.} There is no novelty claim: the proposition derives the observed-law functional from the two declared causal assumptions. \textit{Consumer.} Every boundary approximation and confidence set inherits its causal interpretation from this population identity.

\begin{theorem}[thm:exact-factorization]\label{thm:exact-factorization}
For every \(\ell\) and every admissible \(\rho,c\), \(\log H(r_{\ell1},r_{\ell2},r_{\ell3})-\log(\rho Q_{\ell})=A_{\ell}B_{\rho,\ell}\), \(B_{\rho,\ell}<0\), and \(\Delta_{\ell}-\Delta^P_{\ell}=\rho Q_{\ell}A_{\ell}C_{\rho,\ell}\) with \(C_{\rho,\ell}<0\). Consequently, \(\Delta_{\ell}=\Delta^P_{\ell}\) if and only if \(c_{\ell1}=c_{\ell3}\) or \(c_{\ell2}=c_{\ell3}\).
\end{theorem}
\begin{proof}
Fix \(1\le\ell\le L\).  The coefficient box, rarity-scale assumption, and positivity of \(\rho\) imply
\[
0<\rho c_{\ell j}\le4\rho\le\frac14<1,
\qquad 1\le j\le4.
\tag{1}
\]
Thus \((\rho c_{\ell1},\rho c_{\ell2},\rho c_{\ell3})\in(0,1)^3\), so the global exact-factorization theorem applies.  Under common rarity, put \(p_j=r_{\ell j}=\rho c_{\ell j}\).  The local coordinates in that theorem satisfy
\[
b=p_3=\rho c_{\ell3},
\qquad
x=\log(p_1/p_3)=\log(c_{\ell1}/c_{\ell3})=x_\ell,
\]
\[
z=\log(p_2/p_3)=\log(c_{\ell2}/c_{\ell3})=z_\ell,
\qquad
A=xz=A_\ell.
\tag{2}
\]
Moreover, the Poisson coefficient and the evaluation aliases in the generalized mixed-coefficient definition give
\[
\frac{p_1p_2}{p_3}=\rho\frac{c_{\ell1}c_{\ell2}}{c_{\ell3}}=\rho Q_\ell,
\qquad
B(p_1,p_2,p_3)=B_{\rho,\ell},
\qquad
C(p_1,p_2,p_3)=C_{\rho,\ell}.
\tag{3}
\]
Substitution of (2)--(3) into the global logarithmic identity and sign conclusions proves
\[
\log H(r_{\ell1},r_{\ell2},r_{\ell3})-\log(\rho Q_\ell)
=A_\ell B_{\rho,\ell},
\qquad
B_{\rho,\ell}<0,
\qquad C_{\rho,\ell}<0.
\tag{4}
\]
By the published identification proposition and common rarity,
\[
\Delta_\ell
=H(r_{\ell1},r_{\ell2},r_{\ell3})-r_{\ell4}
=H(r_{\ell1},r_{\ell2},r_{\ell3})-\rho c_{\ell4},
\tag{5}
\]
whereas the definition of the Poisson approximation gives
\[
\Delta^P_\ell=\rho Q_\ell-\rho c_{\ell4}.
\tag{6}
\]
Subtracting (6) from (5) cancels the treated-post term.  The global difference identity and (2)--(3) therefore yield
\[
\Delta_\ell-\Delta^P_\ell
=\rho Q_\ell A_\ell C_{\rho,\ell}.
\tag{7}
\]
Since \(\rho Q_\ell>0\) and \(C_{\rho,\ell}<0\), the right-hand side of (7) vanishes exactly when \(A_\ell=0\).  By (2),
\[
A_\ell=0
\quad\Longleftrightarrow\quad
c_{\ell1}=c_{\ell3}\ \text{or}\ c_{\ell2}=c_{\ell3}.
\]
This proves the existing theorem as the explicit comparable-rarity corollary of the global pointwise result.
\end{proof}
\par\smallskip\noindent \textit{Justification.} This explicit comparable-rarity corollary connects the global population geometry to the panel's \(\rho,c\) notation and causal DTT. \textit{Gap.} Its novel mathematical content is supplied by the global signed-rectangle theorem; the remaining step is the exact rare-panel specialization. \textit{Consumer.} All Gaussian-frontier and common-cutoff descendants use this cancellation-safe rare-panel form.

\begin{theorem}[thm:pairwise-gaussian-frontier]\label{thm:pairwise-gaussian-frontier}
If \(\tau_n=n\rho_n\to\infty\) and \(\Lambda_{\ell n}\to\lambda_{\ell}\in\mathbb R\), then jointly over \(1\le\ell\le L\), \((\widehat\Delta^E_{\ell}-\Delta_{\ell})/\sqrt{\rho_nV_{\ell}/n}\Rightarrow N(0,1)\) and \((\widehat\Delta^P_{\ell}-\Delta_{\ell})/\sqrt{\rho_nV_{\ell}/n}\Rightarrow N(Q_{\ell}\lambda_{\ell}/\sqrt{V_{\ell}},1)\), with independent coordinates. Moreover, the exact-link and Poisson-ratio estimators are first-order equivalent for pair \(\ell\) if and only if \(\Lambda_{\ell n}\to0\).
\end{theorem}
\begin{proof}
Fix \(\ell\), initially suppress that index, and write
\[
 n_j=n_{\ell j},\quad r_{j,n}=\rho_nc_{\ell j,n},\quad
 Q_n=\frac{c_{1,n}c_{2,n}}{c_{3,n}},\quad
 a_n=\sqrt{\frac{\rho_n}{n}}.
\tag{1}
\]
The coefficient-limit and coefficient-box assumptions imply \(c_{j,n}\to c_j\in[1/2,4]\). The role-share limits and bounds imply \(n_j/n\to\pi_j\in[\underline\pi,\overline\pi]\), so every \(\pi_j\) is strictly positive. The rare-limit assumption gives \(\rho_n\to0\), and the hypothesis \(n\rho_n\to\infty\) will justify every comparison with the stochastic scale \(a_n\).
The rarity-scale assumption gives \(0<\rho_n\leq1/16\); together with common rarity and the coefficient box, it yields \(0<r_{j,n}\leq1/4\), so all Bernoulli probabilities and normal quantiles used below are in their declared domains.

The add-half estimator satisfies the exact identity
\[
 \widetilde r_j-r_{j,n}
 =\frac{K_j-n_jr_{j,n}}{n_j+1}
 +\frac{1/2-r_{j,n}}{n_j+1}.
\tag{2}
\]
We first establish the joint cell central limit theorem directly from the declared sampling law. For fixed real numbers \(t_{\ell j}\), put
\[
 h_{ell j,n}=t_{\ell j}\frac{\sqrt{n/\rho_n}}{n_{\ell j}+1}.
\]
The positive role-share lower bound and \(n\rho_n\to\infty\) give \(\max_{ell,j}|h_{ell j,n}|=O((n\rho_n)^{-1/2})=o(1)\). If \(E\sim\operatorname{Bernoulli}(r)\), Taylor's formula for \(e^{ih(E-r)}\), followed by \(\mathbb E(E-r)=0\), gives uniformly for \(|h|\leq1\),
\[
 \mathbb E e^{ih(E-r)}
 =1-\frac{h^2}{2}r(1-r)+R(h,r),
 \qquad |R(h,r)|\leq\frac{|h|^3}{3}r.
\tag{3}
\]
The common-rarity and coefficient-box assumptions bound \(r_{\ell j,n}\leq4\rho_n\). Using \(\log(1+w)=w+O(|w|^2)\) for \(|w|\leq1/2\), equations (3), Bernoulli role sampling, and mutual disjoint-role independence therefore give
\[
 \begin{aligned}
 &\log\mathbb E\exp\!\left{i\sum_{\ell=1}^L\sum_{j=1}^4
 t_{\ell j}\sqrt{\frac n{\rho_n}}
 \frac{K_{ell j,n}-n_{\ell j}r_{\ell j,n}}{n_{\ell j}+1}\right}\\
 &\quad=-\frac12\sum_{\ell=1}^L\sum_{j=1}^4t_{\ell j}^2
 \frac n{\rho_n}\frac{n_{\ell j}r_{\ell j,n}(1-r_{\ell j,n})}{(n_{\ell j}+1)^2}+o(1).
 \end{aligned}
\tag{4}
\]
Indeed, the total cubic remainder for any fixed cell is
\[
 n_{\ell j}|h_{ell j,n}|^3r_{\ell j,n}
 =O((n\rho_n)^{-1/2})=o(1),
\]
and the total quadratic logarithm remainder is \(O(n^{-1})\). The coefficient and role-share limits, common rarity, and \(\rho_n\to0\) imply
\[
 \frac n{\rho_n}\frac{n_{\ell j}r_{\ell j,n}(1-r_{\ell j,n})}{(n_{\ell j}+1)^2}
 \longrightarrow\frac{c_{\ell j}}{\pi_{ell j}}.
\tag{5}
\]
The characteristic function in (4) hence converges to that of a centered \(4L\)-variate normal vector with diagonal covariance having entries \(c_{\ell j}/\pi_{ell j}\). The deterministic term in (2), after multiplication by \(a_n^{-1}=\sqrt{n/\rho_n}\), is \(O((n\rho_n)^{-1/2})=o(1)\). Thus
\[
 a_n^{-1}(\widetilde r-r)
 \Rightarrow N\!\left(0,\operatorname{diag}_{\ell,j}
 \left\{\frac{c_{\ell j}}{\pi_{ell j}}\right\}\right).
\tag{6}
\]
This proves joint convergence and independence across the fixed sixty pair--role cells without importing an interior binomial theorem.

Define the observed-law Poisson functional
\[
 P(p_1,p_2,p_3,p_4)=\frac{p_1p_2}{p_3}-p_4.
\]
At \(r_n=\rho_n(c_{1,n},c_{2,n},c_{3,n},c_{4,n})\), its gradient is
\[
 \nabla P(r_n)=\left(
 \frac{c_{2,n}}{c_{3,n}},
 \frac{c_{1,n}}{c_{3,n}},
 -\frac{c_{1,n}c_{2,n}}{c_{3,n}^2},-1\right).
\tag{7}
\]
All divisions are valid because the coefficient box and positive rarity imply \(r_{3,n}\geq\rho_n/2>0\). From (2), Bernoulli variances, the role-share lower bound, and the coefficient box,
\[
 \max_{\ell,j}\left|\frac{\widetilde r_{\ell j}}{r_{\ell j,n}}-1\right|
 \longrightarrow_P0.
\tag{8}
\]
For example, the squared stochastic part divided by \(r_{\ell j,n}^2\) has expectation \(O((n\rho_n)^{-1})\), the squared add-half part has order \(O((n\rho_n)^{-2})\), and there are only \(4L=60\) cells. Hence the event
\[
 \mathcal E_n=\left\{
 \max_{\ell,j}|\widetilde r_{\ell j}/r_{\ell j,n}-1|\leq1/2
 \right\}
\tag{9}
\]
has probability tending to one. Every first-three-coordinate point on a segment joining \(r_{\ell n}\) to \(\widetilde r_\ell\), on \(\mathcal E_n\), is \(\rho_nd\) for some \(d\in[1/4,6]^3\). Direct differentiation of \(P\) bounds its Hessian on these segments by \(K_P/\rho_n\), where \(K_P<\infty\) is deterministic. Taylor's formula, (6), and \(n\rho_n\to\infty\) now give
\[
 P(\widetilde r_\ell)-P(r_{\ell n})
 =\nabla P(r_{\ell n})^{\!\top}(\widetilde r_\ell-r_{\ell n})
 +O_P(n^{-1})
 =\nabla P(r_{\ell n})^{\!\top}(\widetilde r_\ell-r_{\ell n})+o_P(a_n).
\tag{10}
\]
The variance of the limiting linear form in (10) is
\[
 \begin{aligned}
 V_\ell
 &=\frac{c_{\ell4}}{\pi_{\ell4}}
 +\left(\frac{c_{\ell1}c_{\ell2}}{c_{\ell3}}\right)^2
 \sum_{j=1}^3\frac1{\pi_{\ell j}c_{\ell j}}\\
 &=\frac{c_{\ell4}}{\pi_{\ell4}}+Q_\ell^2
 \sum_{j=1}^3\frac1{\pi_{\ell j}c_{\ell j}}>0.
 \end{aligned}
\tag{11}
\]
Strict positivity follows already from the first summand. Equations (6)--(11) and deterministic coefficient convergence yield, jointly across pairs,
\[
 \frac{\widehat\Delta^P_\ell-\Delta^P_{\ell n}}
 {\sqrt{\rho_nV_\ell/n}}\Rightarrow N(0,1),
\tag{12}
\]
with independent coordinates.

We next transfer (12) to the exact normal-link functional. Invoke the proved fixed-compact Gaussian-tail transfer lemma with the deterministic interval \([d,D]=[1/4,6]\), rarity \(\varepsilon=\rho_n\), and cutoff \(v=u_n=-q(\rho_n)\). The rare-limit assumption makes its fixed cutoff applicable for all sufficiently large \(n\). Uniformly at the true first-three-coordinate vector and at every point of the random Taylor segments on \(\mathcal E_n\), that lemma gives
\[
 \nabla\{H(p_1,p_2,p_3)-p_4\}-\nabla P(p_1,p_2,p_3,p_4)
 =O(u_n^{-2}),
 \qquad
 \max_{j,k\leq3}|\partial_{jk}H|=O(\rho_n^{-1}).
\tag{13}
\]
The fourth-coordinate derivative is exactly \(-1\) for both functionals, and every second derivative involving that coordinate is zero. A second application of Taylor's formula, using (6), (10), and (13), proves
\[
 \begin{aligned}
 &\{H(\widetilde r_{\ell1},\widetilde r_{\ell2},\widetilde r_{\ell3})
 -\widetilde r_{\ell4}-[H(r_{\ell1,n},r_{\ell2,n},r_{\ell3,n})-r_{\ell4,n}]\}\\
 &\quad-\{P(\widetilde r_\ell)-P(r_{\ell n})\}=o_P(a_n).
 \end{aligned}
\tag{14}
\]
Indeed, the gradient-difference term is \(O(u_n^{-2})O_P(a_n)=o_P(a_n)\), while each quadratic remainder is \(O_P(\rho_n^{-1}a_n^2)=O_P(n^{-1})=o_P(a_n)\).

The published-identification proposition, which uses normal-link parallel trends and no anticipation, separately proves the causal equality
\[
 \Delta_{\ell n}=H(r_{\ell1,n},r_{\ell2,n},r_{\ell3,n})-r_{\ell4,n}.
\tag{15}
\]
Thus (15), rather than the definition of either estimator, connects the observed-law functional to the causal target. Combining (12), (14), and (15) proves
\[
 \frac{\widehat\Delta^E_\ell-\Delta_{\ell n}}
 {\sqrt{\rho_nV_\ell/n}}\Rightarrow N(0,1)
\tag{16}
\]
jointly with independent coordinates across \(1\leq\ell\leq L\).

It remains to compute the Poisson centering shift. The exact factorization gives
\[
 \Delta_{\ell n}-\Delta^P_{\ell n}
 =\rho_nQ_{\ell n}A_{\ell n}C_{\rho,\ell n}.
\tag{17}
\]
The uniform common-cutoff theorem applies on the coefficient box, including both exact cancellation planes, and yields
\[
 C_{\rho,\ell n}=-u_n^{-2}\{1+O(u_n^{-2})\}
\tag{18}
\]
uniformly over the fixed pairs. Therefore, without dividing by \(A_{\ell n}\),
\[
 \begin{aligned}
 \frac{\Delta^P_{\ell n}-\Delta_{\ell n}}
 {\sqrt{\rho_nV_\ell/n}}
 &=\frac{Q_{\ell n}}{\sqrt{V_\ell}}
 \frac{A_{\ell n}\sqrt{n\rho_n}}{u_n^2}
 \{1+O(u_n^{-2})\}\\
 &=\frac{Q_\ell}{\sqrt{V_\ell}}\lambda_\ell+o(1),
 \end{aligned}
\tag{19}
\]
because \(Q_{\ell n}\to Q_\ell>0\), \(V_\ell>0\), and \(\Lambda_{\ell n}\to\lambda_\ell\). Adding (19) to (12) proves
\[
 \frac{\widehat\Delta^P_\ell-\Delta_{\ell n}}
 {\sqrt{\rho_nV_\ell/n}}
 \Rightarrow N\!\left(\frac{Q_\ell\lambda_\ell}{\sqrt{V_\ell}},1\right).
\tag{20}
\]
All Taylor remainders above are uniform over the fixed \(L=15\) pairs. Together with the joint characteristic-function calculation in (4), this proves the asserted joint limits; the coordinates indexed by distinct pairs are independent in the limit. The two estimators for a given pair share the same leading Gaussian noise, as equation (14) shows.

For completeness, define first-order equivalence for pair \(\ell\) to mean
\[
 \frac{\widehat\Delta^P_\ell-\widehat\Delta^E_\ell}
 {\sqrt{\rho_nV_\ell/n}}\longrightarrow_P0.
\]
Equations (14), (17), and (18) give
\[
 \frac{\widehat\Delta^P_\ell-\widehat\Delta^E_\ell}
 {\sqrt{\rho_nV_\ell/n}}
 =\frac{Q_{\ell n}}{\sqrt{V_\ell}}\Lambda_{\ell n}
 \{1+O(u_n^{-2})\}+o_P(1).
\tag{21}
\]
The deterministic multiplier converges to \(Q_\ell/\sqrt{V_\ell}\in(0,\infty)\). Hence \(\Lambda_{\ell n}\to0\) implies first-order equivalence. Conversely, if the left side of (21) converges to zero in probability but \(\Lambda_{\ell n}\not\to0\), then a subsequence has its deterministic leading term bounded away from zero, contradicting the \(o_P(1)\) remainder in (21). Thus first-order equivalence holds if and only if \(\Lambda_{\ell n}\to0\), and exact cancellation planes are included because (17)--(21) never divide by \(A_{\ell n}\).
\end{proof}
\par\smallskip\noindent \textit{Justification.} The result turns link disagreement into a drifting noncentrality parameter and now derives the required fixed-compact normal-tail transfer on the estimator's entire high-probability neighborhood. \textit{Gap.} Djuazon and Tsyawo give interior Gaussian weak convergence, while Chen and Fang study degeneracy without vanishing base cell probabilities; neither result supplies the compact-uniform rare-tail value, gradient, and Hessian transfer used here. \textit{Consumer.} A distributional-DiD analyst can determine pair by pair whether the Poisson shortcut changes first-order inference, including on and arbitrarily near either exact cancellation plane.

\begin{theorem}[thm:aggregate-gaussian-frontier]\label{thm:aggregate-gaussian-frontier}
If \(\tau_n\to\infty\) and \(b_n\to b\in\mathbb R\), then \(\{L^{-1}\sum_{\ell=1}^{L}\widehat\Delta^P_{\ell}-\overline\Delta\}/\sqrt{\rho_n\overline V/n}\Rightarrow N(b,1)\), while replacing \(\widehat\Delta^P_{\ell}\) by \(\widehat\Delta^E_{\ell}\) gives \(N(0,1)\). Aggregate first-order equivalence holds if and only if \(b_n\to0\), even when individual pairwise shifts do not vanish.
\end{theorem}
\begin{proof}
Put
\[
 \overline{\widehat\Delta}^{P}_{n}
   =L^{-1}\sum_{\ell=1}^{L}\widehat\Delta^{P}_{\ell n},
 \qquad
 \overline{\widehat\Delta}^{E}_{n}
   =L^{-1}\sum_{\ell=1}^{L}\widehat\Delta^{E}_{\ell n},
 \qquad
 \overline\Delta^{P}_{n}
   =L^{-1}\sum_{\ell=1}^{L}\Delta^{P}_{\ell n},
\]
and let \(s_n=\sqrt{\rho_n/n}\).  We first record the centered consequence of
\(\text{thm:pairwise-gaussian-frontier}\) that is needed below.  Its
characteristic-function and Taylor argument is independent of the convergence
of the drifting indices: under \(\tau_n=n\rho_n\to\infty\), it gives, jointly
over the fixed set \(1\leq\ell\leq L=15\),
\[
 \left(
   \frac{\widehat\Delta^{P}_{\ell n}-\Delta^{P}_{\ell n}}{s_n}
 \right)_{\ell=1}^{L}
 \Rightarrow (Z_\ell)_{\ell=1}^{L},
 \qquad
 \left(
   \frac{\widehat\Delta^{E}_{\ell n}-\Delta_{\ell n}}{s_n}
 \right)_{\ell=1}^{L}
 \Rightarrow (Z_\ell)_{\ell=1}^{L},
 \tag{1}
\]
where the two displays are driven by the same first-order linearized count
vector, the variables \(Z_1,\ldots,Z_L\) are independent, and
\(Z_\ell\sim N(0,V_\ell)\).

For completeness, the absence of a hidden restriction on
\(\Lambda_{\ell n}\) in this centered assertion can be seen directly from the
linearization used in that theorem.  Bernoulli role sampling, common rarity,
the role-share limits, and \(n\rho_n\to\infty\) imply
\[
 \sqrt{\frac n{\rho_n}}
   (\widetilde r_{\ell j}-r_{\ell j})
 =\frac{n}{n_{\ell j}+1}
   \frac{K_{\ell j}-n_{\ell j}r_{\ell j}}{\sqrt{n\rho_n}}
   +o_{P}(1)
 \Rightarrow N\!\left(0,\frac{c_{\ell j}}{\pi_{\ell j}}\right).
 \tag{2}
\]
Indeed, the omitted add-half term is bounded in absolute value by a constant
multiple of \((n\rho_n)^{-1/2}\), using the positive lower role-share bound.
The triangular-array Lindeberg condition follows because every centered
Bernoulli summand, after division by \(\sqrt{n\rho_n}\), is bounded by
\((n\rho_n)^{-1/2}\to0\).  Disjoint-role independence makes all coordinates in
(2) independent.  At the true cell probabilities, the gradient of
\(P^{\rm P}(p_1,p_2,p_3)-p_4\) is
\[
 g_{\ell n}^{P}
 =\left(
   \frac{c_{\ell2,n}}{c_{\ell3,n}},
   \frac{c_{\ell1,n}}{c_{\ell3,n}},
   -\frac{c_{\ell1,n}c_{\ell2,n}}{c_{\ell3,n}^{2}},
   -1
 \right).
 \tag{3}
\]
The coefficient box bounds every denominator in (3) away from zero.  The
Poisson-map Hessian is \(O(\rho_n^{-1})\) on an event of probability tending
to one, because (2) and \(n\rho_n\to\infty\) imply
\(\widetilde r_{\ell3}/r_{\ell3}\to1\) in probability.  Hence its Taylor
remainder is
\[
 O_P\!\left(\rho_n^{-1}\frac{\rho_n}{n}\right)
   =O_P(n^{-1})=o_P(s_n).
 \tag{4}
\]
The fixed-compact Gaussian-tail transfer lemma gives, uniformly on the
coefficient box, that the exact-link gradient differs from (3) by
\(O(u_n^{-2})=o(1)\), and that its Hessian is also
\(O(\rho_n^{-1})\).  Thus (4) applies to the exact-link map as well.  Combining
(2)--(4), the limiting variance for pair \(\ell\) is
\[
 \begin{aligned}
 &\frac{c_{\ell4}}{\pi_{\ell4}}
 +\left(\frac{c_{\ell2}}{c_{\ell3}}\right)^2
    \frac{c_{\ell1}}{\pi_{\ell1}}
 +\left(\frac{c_{\ell1}}{c_{\ell3}}\right)^2
    \frac{c_{\ell2}}{\pi_{\ell2}}
 +\left(\frac{c_{\ell1}c_{\ell2}}{c_{\ell3}^{2}}\right)^2
    \frac{c_{\ell3}}{\pi_{\ell3}} \\
 &\hspace{30mm}
 =\frac{c_{\ell4}}{\pi_{\ell4}}
   +Q_\ell^2\sum_{j=1}^{3}\frac{1}{\pi_{\ell j}c_{\ell j}}
 =V_\ell.
 \end{aligned}
 \tag{5}
\]
Here coefficient and role-share limits justify passage to the displayed
limits, while the rare limit implies \(u_n\to\infty\).  This proves (1)
without imposing convergence or boundedness of any \(\Lambda_{\ell n}\).
Also, every \(c_{\ell j}\) and \(\pi_{\ell j}\) is positive, so (5) gives
\(V_\ell>0\) and therefore
\[
 \overline V=L^{-2}\sum_{\ell=1}^{L}V_\ell>0.
 \tag{6}
\]

Apply the continuous mapping theorem to the fixed linear average in (1).
Independence across pairs and (6) yield
\[
 \frac{\overline{\widehat\Delta}^{P}_{n}-\overline\Delta^{P}_{n}}
      {\sqrt{\rho_n\overline V/n}}
 \Rightarrow N(0,1),
 \qquad
 \frac{\overline{\widehat\Delta}^{E}_{n}-\overline\Delta_n}
      {\sqrt{\rho_n\overline V/n}}
 \Rightarrow N(0,1).
 \tag{7}
\]
The second centering is the aggregate causal estimand because
\(\text{prop:published-identification-baseline}\) identifies
\(\Delta_{\ell n}=H(r_{\ell1,n},r_{\ell2,n},r_{\ell3,n})-r_{\ell4,n}\)
for every pair before averaging.  By \(\text{def:frontier-indices}\), the
remaining centering difference is the exact identity
\[
 \frac{\overline\Delta^{P}_{n}-\overline\Delta_n}
      {\sqrt{\rho_n\overline V/n}}
 =\frac{\sqrt{n/\rho_n}}{L\sqrt{\overline V}}
    \sum_{\ell=1}^{L}(\Delta^{P}_{\ell n}-\Delta_{\ell n})
 =b_n.
 \tag{8}
\]
Consequently, if \(b_n\to b\in\mathbb R\), adding (8) to the first convergence
in (7) and applying Slutsky's theorem proves
\[
 \frac{\overline{\widehat\Delta}^{P}_{n}-\overline\Delta_n}
      {\sqrt{\rho_n\overline V/n}}
 \Rightarrow N(b,1).
 \tag{9}
\]
The second convergence in (7) is the asserted \(N(0,1)\) limit for the
exact-link aggregate.

It remains to prove the equivalence criterion.  The common first-order
linearization in (1), together with the uniform gradient comparison and
Taylor remainder in (4), gives
\[
 \frac{
   (\overline{\widehat\Delta}^{P}_{n}-\overline\Delta^{P}_{n})
   -(\overline{\widehat\Delta}^{E}_{n}-\overline\Delta_n)}
      {\sqrt{\rho_n\overline V/n}}
 =o_P(1).
 \tag{10}
\]
Combining (8) and (10),
\[
 \frac{\overline{\widehat\Delta}^{P}_{n}
       -\overline{\widehat\Delta}^{E}_{n}}
      {\sqrt{\rho_n\overline V/n}}
 =b_n+o_P(1).
 \tag{11}
\]
Thus aggregate first-order equivalence, meaning convergence of the left-hand
side of (11) to zero in probability, holds when \(b_n\to0\).  Conversely, if
that left-hand side converges to zero in probability, subtraction of the
\(o_P(1)\) term in (11) shows that the deterministic sequence \(b_n\) converges
to zero.  This proves the if-and-only-if claim.

Finally, the criterion genuinely allows cross-pair cancellation.  Take two
pairs with constant coefficient triples
\[
 (c_1,c_2,c_3)=(2,2,1),
 \qquad
 (c_1,c_2,c_3)=(1/2,4,1),
 \tag{12}
\]
and put each of the other pairs on an exact cancellation plane.  Both triples
belong to the coefficient box, and direct calculation gives
\[
 Q_1A_1=4(\log2)^2,
 \qquad
 Q_2A_2=-4(\log2)^2.
 \tag{13}
\]
By exact factorization and the uniform common-cutoff expansion,
\[
 \Delta^{P}_{\ell n}-\Delta_{\ell n}
 =-\rho_nQ_\ell A_\ell C_{\rho,\ell}
 =\frac{\rho_nQ_\ell A_\ell}{u_n^2}
   +O\!\left(\frac{\rho_n|Q_\ell A_\ell|}{u_n^4}\right),
 \tag{14}
\]
uniformly for these pairs, including the cancellation planes.  Equations
(13)--(14) cancel the two leading aggregate shifts exactly, and the remaining
pairs contribute zero, so
\[
 \sum_{\ell=1}^{L}(\Delta^{P}_{\ell n}-\Delta_{\ell n})
 =O(\rho_nu_n^{-4}).
 \tag{15}
\]
To display a valid drifting sequence, choose \(\rho_k\downarrow0\), set
\(u_k=-q(\rho_k)\), and choose integer sample scales
\(n_k=\lfloor u_k^6/\rho_k\rfloor\), with role sizes having fixed positive
limiting shares.  Then
\(\sqrt{n_k\rho_k}=u_k^3\{1+o(1)\}\), and (8), (15), and the positive finite
limit in (6) imply \(b_{n_k}=O(u_k^{-1})\to0\).  On the other hand, for each
of the two pairs in (12),
\[
 |\Lambda_{\ell n_k}|
 =|A_\ell|\frac{\sqrt{n_k\rho_k}}{u_k^2}
 =|A_\ell|u_k\{1+o(1)\}\longrightarrow\infty.
 \tag{16}
\]
Hence aggregate first-order equivalence can hold although the contributing
pairwise shifts do not vanish, as claimed.
\end{proof}
\par\smallskip\noindent \textit{Justification.} The exact signed aggregate criterion preserves cross-pair cancellation that a maximum pairwise bound would erase. \textit{Gap.} The published fifteen-pair aggregation assumes interior regularity and has no rare-link noncentrality frontier. \textit{Consumer.} The fixed fifteen-comparison crime aggregate can be assessed without requiring every individual comparison to be safely away from link disagreement.

\begin{theorem}[thm:product-poisson-transition]\label{thm:product-poisson-transition}
If \(\tau_n=n\rho_n\to\tau\in[0,\infty)\), then the sixty-vector \(K\Rightarrow Z\), and jointly \(n\widehat\Delta^P_{\ell}\Rightarrow T_{\ell}\), \(n\widehat\Delta^E_{\ell}\Rightarrow T_{\ell}\), and \(n\Delta_{\ell}\to\tau(Q_{\ell}-c_{\ell4})\). Hence \(nL^{-1}\sum_{\ell=1}^{L}\widehat\Delta^E_{\ell}\Rightarrow L^{-1}\sum_{\ell=1}^{L}T_{\ell}\); at \(\tau=0\), this limit equals \(L^{-1}\sum_{\ell=1}^{L}\{\pi_{\ell3}/(2\pi_{\ell1}\pi_{\ell2})-1/(2\pi_{\ell4})\}\), which need not vanish under unequal shares.
\end{theorem}
\begin{proof}
Write \(N_{\ell j,n}=n_{\ell j}\), \(p_{\ell j,n}=r_{\ell j,n}\), and
\[
\lambda_{\ell j}=\pi_{\ell j}\tau c_{\ell j}.
\]
By \(\texttt{ass:bernoulli-role-sampling}\), \(K_{\ell j,n}\) has the binomial law with parameters \(N_{\ell j,n}\) and \(p_{\ell j,n}\).  By \(\texttt{ass:common-rarity}\), \(\texttt{ass:role-share-limits}\), \(\texttt{ass:coefficient-limits}\), and the displayed premise \(n\rho_n\to\tau\),
\[
N_{\ell j,n}p_{\ell j,n}
=\frac{N_{\ell j,n}}n(n\rho_n)c_{\ell j,n}
\longrightarrow\lambda_{\ell j}.
\tag{17}
\]
The rarity limit and coefficient limits imply \(p_{\ell j,n}=\rho_nc_{\ell j,n}\to0\), and therefore
\[
N_{\ell j,n}p_{\ell j,n}^2
=(N_{\ell j,n}p_{\ell j,n})p_{\ell j,n}\longrightarrow0.
\tag{18}
\]
Moreover, \(\texttt{ass:role-share-bounds}\) and the share limits imply
\(\underline\pi\leq\pi_{\ell j}\leq\overline\pi\); in particular, every \(\pi_{\ell j}\) is positive and \(N_{\ell j,n}\to\infty\).

For every fixed nonnegative integer \(k\), the exact binomial mass is
\[
\Pr(K_{\ell j,n}=k)
=\frac1{k!}\prod_{h=0}^{k-1}(N_{\ell j,n}-h)p_{\ell j,n}
\,(1-p_{\ell j,n})^{N_{\ell j,n}}(1-p_{\ell j,n})^{-k}.
\tag{19}
\]
Equations (17)--(18) give
\[
N_{\ell j,n}\log(1-p_{\ell j,n})
=-N_{\ell j,n}p_{\ell j,n}+O(N_{\ell j,n}p_{\ell j,n}^2)
\longrightarrow-\lambda_{\ell j}.
\tag{20}
\]
For \(k=0\), the product in (19) is empty; for \(k\geq1\), each of its \(k\) factors converges to \(\lambda_{\ell j}\).  Thus (19)--(20) prove
\[
\Pr(K_{\ell j,n}=k)
\longrightarrow e^{-\lambda_{\ell j}}\frac{\lambda_{\ell j}^k}{k!},
\tag{21}
\]
with the same conclusion when \(\lambda_{\ell j}=0\).

The disjoint-role independence assumption makes the sixty counts mutually independent.  Hence their joint mass at each point of \(\mathbb N^{4L}\) converges, by (21), to the product of the Poisson masses in \(\texttt{def:bounded-count-limit}\).  In addition, (17) and the fixed dimension give a finite constant \(C_0\) such that, for all sufficiently large \(n\),
\[
\mathbb E\!\left[\sum_{\ell=1}^{L}\sum_{j=1}^{4}K_{\ell j,n}\right]
=\sum_{\ell=1}^{L}\sum_{j=1}^{4}N_{\ell j,n}p_{\ell j,n}
\leq C_0.
\tag{22}
\]
Therefore
\[
\Pr\!\left(\max_{\ell,j}K_{\ell j,n}>M\right)
\leq\frac{C_0}{M+1}.
\tag{23}
\]
Pointwise convergence of the joint masses on the finite cube \(\{0,\ldots,M\}^{4L}\), followed by (23) and then \(M\to\infty\), proves directly that
\[
K\Rightarrow Z,
\tag{24}
\]
where the coordinates of \(Z\) are independent and have the stated Poisson means.  Indeed, the same Markov bound applies to \(Z\), because \(\mathbb E\sum_{\ell,j}Z_{\ell j}=\sum_{\ell,j}\lambda_{\ell j}<\infty\); hence the omitted mass outside the finite cube tends to zero for both sides.

Define the positive rescaled add-half cells
\[
a_{\ell j,n}=n\widetilde r_{\ell j}
=\frac{K_{\ell j,n}+1/2}{(N_{\ell j,n}+1)/n}.
\tag{25}
\]
The deterministic denominators in (25) converge to \(\pi_{\ell j}>0\).  On every fixed finite count cube, the map from the count vector to the vector in (25) therefore converges uniformly to
\[
k\longmapsto\left(\frac{k_{\ell j}+1/2}{\pi_{\ell j}}\right)_{\ell,j}.
\]
Combining this finite-cube uniform convergence with (23)--(24) proves, without any unbounded mapping step,
\[
(a_{\ell j,n})_{\ell,j}\Rightarrow
(W_{\ell j})_{\ell,j},
\qquad W_{\ell j}=\frac{Z_{\ell j}+1/2}{\pi_{\ell j}}.
\tag{26}
\]
Every coordinate in (25)--(26) is strictly positive, including when its count is zero.  The same finite-cube argument applied to the rational map gives jointly over \(1\leq\ell\leq L\),
\[
n\widehat\Delta^P_{\ell}
=\frac{a_{\ell1,n}a_{\ell2,n}}{a_{\ell3,n}}-a_{\ell4,n}
\Rightarrow
\frac{W_{\ell1}W_{\ell2}}{W_{\ell3}}-W_{\ell4}
=T_{\ell}.
\tag{27}
\]
The divisions in (27) are valid because the add-half numerator is at least \(1/2\).

We next compare the exact and multiplicative links.  Fix an integer \(M\geq0\), and define
\[
\mathcal F_{n,M}=\left\{\max_{\ell,j}K_{\ell j,n}\leq M\right\},
\qquad
d_M=\frac1{2(\overline\pi+1)},
\qquad
D_M=\frac{M+1/2}{\underline\pi}.
\tag{28}
\]
On \(\mathcal F_{n,M}\), the share bounds and (25) give, for every cell and every \(n\),
\[
0<d_M\leq a_{\ell j,n}\leq D_M<\infty.
\tag{29}
\]
Indeed, \(N_{\ell j,n}+1\leq n(\overline\pi+1)\) gives the lower bound, and \(N_{\ell j,n}+1\geq n\underline\pi\) gives the upper bound.  Apply \(\texttt{lem:fixed-compact-gaussian-tail-transfer}\) to this deterministic interval \([d_M,D_M]\), with \(\varepsilon=1/n\).  If \(v_n=-q(1/n)\), then \(v_n\to\infty\), and there is a finite constant \(C_M\), independent of \(n\), the pair, and the realized count vector in \(\mathcal F_{n,M}\), such that
\[
\max_{1\leq\ell\leq L}
\left|
nH(a_{\ell1,n}/n,a_{\ell2,n}/n,a_{\ell3,n}/n)
-\frac{a_{\ell1,n}a_{\ell2,n}}{a_{\ell3,n}}
\right|
\leq\frac{C_M}{v_n^2}
\quad\hbox{on }\mathcal F_{n,M}.
\tag{30}
\]
The arguments of \(H\) are exactly the add-half probabilities and hence lie in \((0,1)\), even for all-zero or all-full count cells.  By (23), for every \(\delta>0\), (30) implies
\[
\limsup_{n\to\infty}
\Pr\!\left(
\max_{1\leq\ell\leq L}
|n\widehat\Delta^E_{\ell}-n\widehat\Delta^P_{\ell}|>\delta
\right)
\leq\frac{C_0}{M+1}.
\tag{31}
\]
Letting \(M\to\infty\) proves
\[
\max_{1\leq\ell\leq L}
|n\widehat\Delta^E_{\ell}-n\widehat\Delta^P_{\ell}|
\longrightarrow_P0.
\tag{32}
\]
Equations (27) and (32), again checked first on finite count cubes and then using (23), prove the asserted joint exact-link statistic limit.

It remains to relate the observed-law functional to the causal target.  The finite collection of positive limits in \(\texttt{ass:coefficient-limits}\) permits fixed constants \(0<d_\star<D_\star<\infty\) such that
\[
c_{\ell j,n}\in[d_\star,D_\star]
\quad\hbox{for every }\ell,j\hbox{ and all sufficiently large }n.
\tag{33}
\]
Apply the fixed-compact transfer lemma to \([d_\star,D_\star]\), now with \(\varepsilon=\rho_n\).  Put \(u_n=-q(\rho_n)\), and define the proof intermediate
\[
Q_{\ell,n}=\frac{c_{\ell1,n}c_{\ell2,n}}{c_{\ell3,n}}.
\]
The rarity limit gives \(u_n\to\infty\), and the lemma yields a finite \(C_\star\) such that
\[
\max_{1\leq\ell\leq L}
\left|
\rho_n^{-1}H(\rho_nc_{\ell1,n},\rho_nc_{\ell2,n},\rho_nc_{\ell3,n})
-Q_{\ell,n}
\right|
\leq\frac{C_\star}{u_n^2}.
\tag{34}
\]
By \(\texttt{prop:published-identification-baseline}\), the causal contrast is identified by the observed-law map,
\[
\Delta_{\ell,n}
=H(r_{\ell1,n},r_{\ell2,n},r_{\ell3,n})-r_{\ell4,n}.
\tag{35}
\]
Substitution of common rarity into (35), followed by (34), gives
\[
n\Delta_{\ell,n}
=n\rho_n\{Q_{\ell,n}-c_{\ell4,n}\}
+O\!\left(\frac{n\rho_n}{u_n^2}\right).
\tag{36}
\]
The sequence \(n\rho_n\) is bounded, \(u_n\to\infty\), and coefficient convergence with \(c_{\ell3}>0\) gives \(Q_{\ell,n}\to Q_{\ell}\).  Hence (36) proves
\[
n\Delta_{\ell,n}\longrightarrow
\tau(Q_{\ell}-c_{\ell4}).
\tag{37}
\]
This establishes the claimed relationship between the causal target and the observed-law transfer; it is not imposed by definition.

Because \(L=15\) is fixed, summing the joint exact-link limits and dividing by \(L\) gives
\[
nL^{-1}\sum_{\ell=1}^{L}\widehat\Delta^E_{\ell}
\Rightarrow L^{-1}\sum_{\ell=1}^{L}T_{\ell}.
\tag{38}
\]
Finally, if \(\tau=0\), then every Poisson mean is zero and every \(Z_{\ell j}=0\) almost surely.  Therefore
\[
W_{\ell j}=\frac1{2\pi_{\ell j}},
\qquad
T_{\ell}
=\frac{\pi_{\ell3}}{2\pi_{\ell1}\pi_{\ell2}}
-\frac1{2\pi_{\ell4}},
\tag{39}
\]
which proves the displayed aggregate limit.  Positivity of all shares validates every denominator.  Unequal shares do not force (39) to be zero: whenever the share bounds permit two distinct positive values \(a\) and \(b\), shares satisfying \(\pi_{\ell1}=\pi_{\ell2}=\pi_{\ell3}=a\) and \(\pi_{\ell4}=b\neq a\) give \(T_{\ell}=1/(2a)-1/(2b)\neq0\).  Thus the all-zero limit is retained by the add-half rule and may exhibit the stated unequal-share displacement.
\end{proof}
\par\smallskip\noindent \textit{Justification.} Bounded expected counts change the experiment itself, so a Gaussian correction cannot describe the sampling law. \textit{Gap.} Marohn supplies general Gaussian--Poisson experiment transitions but not the smoothed nonlinear DiD statistic or its unequal-share zero-intensity limit. \textit{Consumer.} Sparse panels with zero cells obtain a defined non-Gaussian reference law rather than an invalid Wald approximation.

\begin{lemma}[lem:global-lipschitz]\label{lem:global-lipschitz}
For every \(p,p'\in\mathcal P_0\), \(|G(p)-G(p')|\le L_{\rm lip}\lVert p-p'\rVert_{\infty}\).
\end{lemma}
\begin{proof}
The set \(\mathcal P_0\) is an intersection of closed linear half-spaces and a box, hence is compact and convex.  If \(p\in\mathcal P_0\) has a zero coordinate, the inequalities \(p_i\le8p_j\) force every coordinate to be zero.  Thus every nonzero point is strictly positive.

At a positive point put \(t=\max(p_1,p_2,p_3)\), \(v=-q(t)\), and write
\[
q(p_i)=-v-d_i,\qquad i=1,2,3.
\]
Since \(\phi(w)\le\phi(0)=(2\pi)^{-1/2}<1/2\),
\[
\Phi(-1/2)=\frac12-\int_0^{1/2}\phi(w)\,dw>\frac14.
\]
Because \(t\le1/4\), one therefore has \(v>1/2\).  Comparable rarity gives \(t/8\le p_i\le t\), and the identity
\[
\log\frac{t}{p_i}=\int_v^{v+d_i}m(-w)\,dw
\]
together with \(m(-w)>w\) gives
\[
0\le d_i\le d:=\frac{\log8}{v}.
\tag{1}
\]
Let \(s=q(p_1)+q(p_2)-q(p_3)\).  The elementary bounds following from (1) are
\[
|s-q(p_i)|\le2d,
\qquad |s+q(p_i)|\le2v+3d.
\]
Since \(\log8<3\) and \(v^{-2}<4\),
\[
\left|\frac{s^2-q(p_i)^2}{2}\right|
\le2\log8+\frac{3(\log8)^2}{v^2}<114.
\tag{2}
\]
At positive points,
\[
\left|\frac{\partial H}{\partial p_i}\right|
=\frac{\phi(s)}{\phi(q(p_i))}
=\exp\left\{\frac{q(p_i)^2-s^2}{2}\right\}<e^{114}<2^{256}.
\]
The last numerical inequality follows from \(\log2>1/2\).  Since the derivative of \(-p_4\) is \(-1\), (2) gives
\[
\|\nabla G(p)\|_1<3\cdot2^{256}+1<2^{258}=L_{\rm lip}
\tag{3}
\]
throughout the positive part of the cone.

We also verify the boundary value needed to integrate (3).  Along any positive sequence tending to the origin, set \(\rho=t/4\) and \(c_i=p_i/\rho\) for \(i=1,2,3\).  Then \(c_i\in[1/2,4]\) and \(\rho\downarrow0\).  The inverse-tail integral identity
\[
\frac{\Phi(-v)}{\phi(v)}=\frac1v\int_0^\infty e^{-a}e^{-a^2/(2v^2)}\,da
\]
and (1) show that the combined normal index equals \(-v+e_v\), where \(|e_v|\le3\log(8)/v\).  Moreover, for \(w\ge1\), restricting the integral to \([0,1]\) gives
\[
\frac{\Phi(-w)}{\phi(w)}\ge
\frac1w\int_0^1e^{-a-a^2/2}\,da=\frac{c_0}{w},
\qquad c_0:=\int_0^1e^{-a-a^2/2}\,da>0,
\]
and hence \(m(-w)\le w/c_0\).  If \(e_v\le0\), monotonicity gives \(\Phi(-v+e_v)\le\Phi(-v)\).  If \(e_v>0\), then for all sufficiently large \(v\),
\[
\log\frac{\Phi(-v+e_v)}{\Phi(-v)}
=\int_{v-e_v}^{v}m(-w)\,dw
\le \frac{ve_v}{c_0}\le\frac{3\log8}{c_0}.
\]
Thus \(H(p_1,p_2,p_3)=O\{\Phi(-v)\}=O(t)=O(\rho)\).  The cone inequalities also give \(p_4=O(t)\), so \(G(p)\to0=G(0)\).  Thus \(G\) is continuous at the only boundary point where differentiability was not asserted.

For positive \(p,p'\), the line segment joining them stays in the positive part of the convex cone, so integration of (3) gives
\[
|G(p)-G(p')|\le L_{\rm lip}\|p-p'\|_\infty.
\]
If one endpoint is the origin, apply this inequality to an interior point of the segment and pass to the origin using the proved continuity.  This establishes the global certificate.
\end{proof}
\par\smallskip\noindent \textit{Justification.} A global certificate transfers continuum coverage through a finite rational grid, including the origin. \textit{Gap.} General exact multinomial-function inversion assumes continuity but does not give a usable global constant for this inverse-normal composition. \textit{Consumer.} Auditable software can bound off-grid error without floating-point sampling of derivatives near zero.

\begin{theorem}[thm:rational-outer-interval]\label{thm:rational-outer-interval}
For rational \(\beta\in(0,1)\) and \(0<\eta\le1/n\), \(\mathsf{Outer}_{\beta,\eta}\) terminates after finitely many rational operations and certified normal-CDF enclosures, returns a rational interval containing \(G(\mathcal R_{\beta}(K))\), and satisfies
\[
\inf_{p\in\mathcal P_0}\Pr_p\{G(p)\in\mathsf{Outer}_{\beta,\eta}(K)\}\ge1-\beta.
\]
There exists a finite constant \(C_{\beta}\), depending only on the fixed role-share bounds, such that
\[
\sup_{p=\rho c}\mathbb E_p[\operatorname{diam}\mathsf{Outer}_{\beta,\eta}(K)]
\le C_{\beta}\{\max(1,\sqrt{n\rho})/n+\eta\}.
\]
The all-zero table retains the origin.  The all-full table produces an empty retained list whenever \(2\sum_{j=1}^{4}n_j\log4>8/\beta+2\log4\).
\end{theorem}
\begin{proof}
Fix one pair, write \(n_j\) for its four role sizes, \(M_n=\sum_jn_j\), and let \(x_j=k_j/n_j\).  Define the polynomial likelihood without combinatorial factors by
\[
L_k(p)=\prod_{j=1}^4p_j^{k_j}(1-p_j)^{n_j-k_j},
\qquad
L_{\rm sat}=L_k(x)>0,
\]
with \(0^0=1\).  The endpoint conventions in the Bernoulli divergence give the exact extended-real identity
\[
D_p(k)=2\{\log L_{\rm sat}-\log L_k(p)\}.
\tag{1}
\]
Put \(t=L_{\rm sat}e^{-4/\beta}\).  Certified rational upper and lower exponential series bounds produce, after finitely many terms, a positive rational \(t_{\rm lo}\) satisfying
\[
t/2\le t_{\rm lo}\le t.
\tag{2}
\]
Choose a positive rational mesh \(h\) with
\[
M_nh\le t_{\rm lo}/4,
\qquad L_{\rm lip}h\le\eta/8,
\tag{3}
\]
and partition \([0,1/4]^4\) into finitely many rational boxes of max-diameter at most \(h\).

For each box \(B\), feasibility of \(B\cap\mathcal P_0\) is decidable using rational arithmetic: enumerate quadruples of active rational boundary hyperplanes, solve every nonsingular rational system, and check its solution against all inequalities.  A nonempty bounded rational polytope has a vertex, and at a vertex the active normals span \(\mathbb R^4\), so this finite enumeration also returns a rational feasible point \(p_B\).  For one coordinate, the absolute derivative of \(p^k(1-p)^{n-k}\) is at most \(n\) on \([0,1]\).  Telescoping over four coordinates therefore gives
\[
|L_k(p)-L_k(p')|\le M_n\|p-p'\|_\infty.
\tag{4}
\]
Thus \(U_B=L_k(p_B)+M_nh\) is a rational upper enclosure on \(B\).  Retain exactly the feasible boxes with \(U_B\ge t_{\rm lo}\).

At a positive rational feasible point, \(G(p_B)\) has a certified rational enclosure of any prescribed positive width.  One completely finite construction is to enumerate rational lower and upper brackets for each inverse normal quantile and certify strict CDF inequalities by rational Taylor bounds for the exponential integral; a bracket of the requested width exists by continuity and strict monotonicity.  Enumerating pairs, rather than bisecting at a possibly exact equality, guarantees termination.  Applying the same certified CDF enclosure to the resulting normal-index interval gives a rational interval \([a_B,b_B]\ni G(p_B)\) of width at most \(\eta/4\).  At \(p_B=0\), use the exact enclosure \([0,0]\).  The global Lipschitz lemma implies
\[
G(B\cap\mathcal P_0)\subset
[a_B-L_{\rm lip}h,b_B+L_{\rm lip}h].
\tag{5}
\]
Take the rational hull of these intervals and round outward on a rational grid of width at most \(\eta/2\); if there is no retained box, return \([0,0]\).  Every enumeration, series refinement, and grid is finite, proving termination and rationality.

If \(p\in\mathcal R_\beta(k)\), then (1) gives \(L_k(p)\ge t\ge t_{\rm lo}\).  Its feasible box satisfies \(U_B\ge L_k(p)\), is retained, and (5) contains \(G(p)\).  Hence the returned interval contains \(G(\mathcal R_\beta(k))\).  An empty retained list therefore implies \(\mathcal R_\beta(k)=\varnothing\).

For coverage, first take positive \(p\in\mathcal P_0\).  The elementary inequality \(\log y\le y-1\) gives
\[
\operatorname{KL}\{\operatorname{Bern}(x),\operatorname{Bern}(p)\}
\le\frac{(x-p)^2}{p(1-p)}.
\]
Therefore, under the four independent binomial marginals,
\[
\mathbb E_pD_p(K)
\le2\sum_{j=1}^4n_j
\frac{\mathbb E_p(K_j/n_j-p_j)^2}{p_j(1-p_j)}=8.
\]
Markov's inequality yields
\[
\Pr_p\{p\notin\mathcal R_\beta(K)\}
=\Pr_p\{D_p(K)>8/\beta\}\le\beta.
\tag{6}
\]
At the origin, \(K=0\) almost surely, the origin's box obeys \(U_B\ge L_k(0)=L_{\rm sat}=1\), and so it is retained.  Thus (6) and the image containment prove the uniform \(1-\beta\) coverage, including the origin.

It remains to prove the expected-diameter order.  If \(B\) is retained, (3)--(4) imply
\[
L_k(p_B)\ge t_{\rm lo}-M_nh\ge3t_{\rm lo}/4.
\]
Every feasible point in that box has likelihood at least \(t_{\rm lo}/2\ge t/4\), and hence, by (1), belongs to
\[
\mathcal R'(k)=\{p\in\mathcal P_0:D_p(k)\le q'\},
\qquad q'=8/\beta+2\log4.
\tag{7}
\]
For Bernoulli laws, Jensen's inequality applied to \(-\log\), followed by \(-\log a\ge1-a\), gives
\[
\operatorname{KL}\{\operatorname{Bern}(x),\operatorname{Bern}(p)\}
\ge(\sqrt x-\sqrt p)^2+(\sqrt{1-x}-\sqrt{1-p})^2.
\]
Consequently every \(p\in\mathcal R'(k)\) satisfies
\[
|\sqrt{p_j}-\sqrt{x_j}|\le d_j:=\sqrt{\frac{q'}{2n_j}}.
\tag{8}
\]
For any two such points, (8) gives
\[
|p_j-p'_j|\le4\sqrt{x_j}d_j+4d_j^2.
\tag{9}
\]
The Lipschitz image bound, (5), and the final outward rounding show that the output diameter is at most a fixed multiple of
\[
L_{\rm lip}\sum_{j=1}^4\{\sqrt{x_j/n_j}+1/n_j\}+\eta.
\]
Since \(\mathbb E_p\sqrt{x_j}\le\sqrt{p_j}\), \(p_j=\rho c_j\le4\rho\), and \(n_j\ge\underline\pi n\), taking expectations proves
\[
\sup_{p=\rho c}\mathbb E_p\operatorname{diam}\mathsf{Outer}_{\beta,\eta}(K)
\le C_\beta\left\{\sqrt{\frac\rho n}+\frac1n+\eta\right\}
\le C_\beta\left\{\frac{\max(1,\sqrt{n\rho})}{n}+\eta\right\}.
\]
Here \(C_\beta<\infty\) depends only on \(\beta\), the fixed share bounds, and the fixed certified Lipschitz constant.

For \(k=0\), the origin argument above proves retention.  If instead \(k_j=n_j\) for every \(j\), then \(L_{\rm sat}=1\) and every feasible \(p\) satisfies \(L_k(p)\le4^{-\sum_jn_j}\).  A retained feasible point would, by (2)--(4), satisfy \(L_k(p_B)\ge t/4\), which is impossible when
\[
2\sum_{j=1}^4n_j\log4>8/\beta+2\log4.
\]
Thus the retained list is empty under the stated all-full condition and the algorithm returns exactly \([0,0]\).
\end{proof}
\par\smallskip\noindent \textit{Justification.} The theorem gives boundary-valid coverage on every count table together with a finite certificate and a rare-regime length guarantee. \textit{Gap.} Sachs, Gabriel, and Fay provide valid continuum inversion and a stochastic approximation, not a deterministic outer certificate with this diameter order. \textit{Consumer.} Small-block distributional-DiD studies can report intervals that remain defined at zero counts and can be independently audited.

\begin{proposition}[prop:simultaneous-aggregate-interval]\label{prop:simultaneous-aggregate-interval}
Suppose \(\alpha\in\mathbb Q\cap(0,1)\), set \(\beta=\alpha/L\), and let \(\eta>0\) be the predeclared rational tolerance.  Then
\[
\Pr\{\Delta_\ell\in I_\ell\text{ for every }1\le\ell\le L\}\ge1-\alpha
\quad\text{and}\quad
\Pr\{\overline\Delta\in\overline I\}\ge1-\alpha.
\]
There exists a finite constant \(C_\alpha\), depending only on \(\alpha\) and the fixed role-share bounds, such that
\[
\mathbb E\{\operatorname{diam}(\overline I)\}
\le C_\alpha\left\{\frac{\max(1,\sqrt{n\rho})}{n}+\eta\right\}.
\]
\end{proposition}
\begin{proof}
Fix \(n\) and \(1\leq \ell\leq L\), and define the observed-law probability vector
\[
p_\ell=(r_{\ell1},r_{\ell2},r_{\ell3},r_{\ell4}).
\]
The common-rarity assumption gives \(p_{\ell j}=\rho c_{\ell j}\).  By the coefficient-box and rarity-scale assumptions,
\[
0<p_{\ell j}\leq4\rho\leq\frac14
\quad\text{and}\quad
\frac{p_{\ell i}}{p_{\ell j}}
=\frac{c_{\ell i}}{c_{\ell j}}
\leq\frac4{1/2}=8
\qquad(i,j\in J).
\tag{1}
\]
Consequently \(p_\ell\in\mathcal P_0\).  The published-identification proposition, which uses normal-link parallel trends and no anticipation, now yields the population identity
\[
\Delta_\ell
=H(r_{\ell1},r_{\ell2},r_{\ell3})-r_{\ell4}
=G(p_\ell).
\tag{2}
\]
Thus (2) proves the relationship between the causal contrast \(\Delta_\ell\) and the observed-law functional \(G(p_\ell)\); it does not define the causal contrast to equal that functional.

Set \(\beta=\alpha/L\).  Since \(L=15\), \(\alpha\in\mathbb Q\cap(0,1)\) implies \(\beta\in\mathbb Q\cap(0,1)\).  The rational outer-interval theorem is stated for \(\eta\leq1/n\).  We first verify that the two conclusions used here continue to hold for every positive rational \(\eta\).

For the fixed pair, abbreviate \(n_j=n_{\ell j}\), \(K_j=K_{\ell j,n}\), \(x_j=K_j/n_j\), and \(M_n=\sum_{j=1}^4n_j\).  Write
\[
L_K(p)=\prod_{j=1}^4p_j^{K_j}(1-p_j)^{n_j-K_j},
\qquad
L_{\rm sat}=L_K(x),
\tag{3}
\]
with \(0^0=1\).  Then \(L_{\rm sat}>0\), and the endpoint conventions in the deviance definition give
\[
D_p(K)=2\{\log L_{\rm sat}-\log L_K(p)\}.
\tag{4}
\]
The rational algorithm chooses \(t_{\rm lo}>0\) such that, on putting \(t=L_{\rm sat}e^{-4/\beta}\),
\[
\frac t2\leq t_{\rm lo}\leq t.
\tag{5}
\]
For every rational \(\eta>0\), density of the positive rationals provides a rational \(h>0\) satisfying
\[
M_nh\leq\frac{t_{\rm lo}}4,
\qquad
L_{\rm lip}h\leq\frac\eta8.
\tag{6}
\]
Hence the grid in the algorithm is finite.  The likelihood is globally Lipschitz on \([0,1]^4\): telescoping first over the \(n_j\) Bernoulli factors in role \(j\), and then over the four roles, gives
\[
|L_K(p)-L_K(p')|
\leq M_n\lVert p-p'\rVert_\infty.
\tag{7}
\]
If \(p\in\mathcal R_\beta(K)\), (4) gives \(L_K(p)\geq t\geq t_{\rm lo}\).  Therefore the rational likelihood upper enclosure of the feasible grid box containing \(p\) meets \(t_{\rm lo}\), so that box is retained.  The global Lipschitz certificate for \(G\), the second inequality in (6), and outward rounding in the rational algorithm then imply
\[
G\{\mathcal R_\beta(K)\}
\subseteq \mathsf{Outer}_{\beta,\eta}(K).
\tag{8}
\]
The finite certified normal-CDF enclosure argument from the rational outer-interval theorem uses only \(\eta>0\), through (6), so it also terminates here.  Thus neither termination nor (8) requires \(\eta\leq1/n\).

Under Bernoulli role sampling and disjoint-role independence, \(K_1,\ldots,K_4\) are independent with \(K_j\sim\operatorname{Binomial}(n_j,p_{\ell j})\).  For \(x\in[0,1]\) and \(p\in(0,1)\), applying \(\log y\leq y-1\) to the two terms of the Bernoulli divergence gives
\[
\operatorname{KL}\{\operatorname{Bern}(x),\operatorname{Bern}(p)\}
\leq\frac{(x-p)^2}{p(1-p)}.
\tag{9}
\]
Since \(0<p_{\ell j}\leq1/4\), division in (9) is valid.  Moreover,
\[
\mathbb E_{p_\ell}(x_j-p_{\ell j})^2
=\frac{p_{\ell j}(1-p_{\ell j})}{n_j}.
\tag{10}
\]
Equations (9)--(10) and the deviance definition therefore yield
\[
\mathbb E_{p_\ell}D_{p_\ell}(K)
\leq2\sum_{j=1}^4n_j
\frac{p_{\ell j}(1-p_{\ell j})/n_j}
{p_{\ell j}(1-p_{\ell j})}
=8.
\tag{11}
\]
Markov's inequality, the definition of \(\mathcal R_\beta\), (8), and (2) imply
\[
\begin{aligned}
\Pr_{p_\ell}\{\Delta_\ell\notin I_\ell\}
&=\Pr_{p_\ell}\{G(p_\ell)\notin I_\ell\}\\
&\leq\Pr_{p_\ell}\{p_\ell\notin\mathcal R_\beta(K)\}\\
&=\Pr_{p_\ell}\{D_{p_\ell}(K)>8/\beta\}\\
&\leq\beta.
\end{aligned}
\tag{12}
\]

We also extend the diameter calculation to arbitrary \(\eta>0\).  Let \(B\) be a retained box and let \(p_B\in B\cap\mathcal P_0\) be its rational feasible point.  By retention, (6), and (7),
\[
L_K(p_B)\geq t_{\rm lo}-M_nh\geq\frac{3t_{\rm lo}}4.
\tag{13}
\]
For every \(p\in B\cap\mathcal P_0\), another application of (6)--(7) gives
\[
L_K(p)\geq L_K(p_B)-M_nh
\geq\frac{t_{\rm lo}}2
\geq\frac t4.
\tag{14}
\]
Combining (4), (5), and (14), every feasible point of every retained box belongs to
\[
\mathcal R'(K)
=\left\{p\in\mathcal P_0:D_p(K)\leq q'\right\},
\qquad
q'=\frac8\beta+2\log4.
\tag{15}
\]
For Bernoulli laws, the elementary Hellinger lower bound
\[
\operatorname{KL}\{\operatorname{Bern}(x),\operatorname{Bern}(p)\}
\geq(\sqrt{x}-\sqrt{p})^2
+(\sqrt{1-x}-\sqrt{1-p})^2
\tag{16}
\]
follows by applying Jensen's inequality to \(-\log\) and then using \(-\log a\geq1-a\).  Since every summand defining \(D_p(K)\) is nonnegative, (15)--(16) imply, for \(p\in\mathcal R'(K)\),
\[
|\sqrt{p_j}-\sqrt{x_j}|
\leq d_j,
\qquad
d_j:=\sqrt{\frac{q'}{2n_j}}.
\tag{17}
\]
For any \(p,p'\in\mathcal R'(K)\), (17) consequently gives
\[
|p_j-p'_j|
\leq2d_j(2\sqrt{x_j}+2d_j)
=4\sqrt{x_j}\sqrt{\frac{q'}{2n_j}}+\frac{2q'}{n_j}.
\tag{18}
\]
The global Lipschitz bound for \(G\), the numerical enclosure and Lipschitz enlargement in the rational algorithm, (6), and its final outward rounding now give a finite constant \(C_{0,\beta}\), independent of \(\ell,n,\rho\), and \(c\), such that
\[
\operatorname{diam}(I_\ell)
\leq C_{0,\beta}
\left\{
\sum_{j=1}^4\left(\sqrt{\frac{x_j}{n_j}}+\frac1{n_j}\right)+\eta
\right\}.
\tag{19}
\]
If no box is retained, the algorithm returns \([0,0]\), and (19) remains valid.  Jensen's inequality, common rarity, the upper coefficient bound, and the lower role-share bound imply
\[
\mathbb E_{p_\ell}\sqrt{x_j}
\leq\sqrt{\mathbb E_{p_\ell}x_j}
=\sqrt{p_{\ell j}}
\leq2\sqrt\rho,
\qquad
n_j\geq\underline\pi n.
\tag{20}
\]
Taking expectations in (19) and using (20), uniformly over the fixed pairs, yields a finite \(C_\beta\), depending only on \(\beta\), the fixed role-share bounds, and the fixed certificate \(L_{\rm lip}\), such that
\[
\mathbb E_{p_\ell}\{\operatorname{diam}(I_\ell)\}
\leq C_\beta\left\{\sqrt{\frac\rho n}+\frac1n+\eta\right\}
\leq C_\beta\left\{\frac{\max(1,\sqrt{n\rho})}{n}+\eta\right\}.
\tag{21}
\]
This proves for every positive rational \(\eta\) the pairwise coverage and diameter facts needed below.

Define \(A_\ell=\{\Delta_\ell\notin I_\ell\}\).  No independence across pairs is needed: the union bound and (12) give
\[
\begin{aligned}
\Pr\{\Delta_\ell\in I_\ell\text{ for every }1\leq\ell\leq L\}
&=1-\Pr\left(\bigcup_{\ell=1}^LA_\ell\right)\\
&\geq1-\sum_{\ell=1}^L\Pr(A_\ell)\\
&\geq1-L\beta
=1-\alpha.
\end{aligned}
\tag{22}
\]
Write \(I_\ell=[L_\ell,U_\ell]\).  On the event in (22),
\[
L^{-1}\sum_{\ell=1}^LL_\ell
\leq L^{-1}\sum_{\ell=1}^L\Delta_\ell
=\overline\Delta
\leq L^{-1}\sum_{\ell=1}^LU_\ell.
\tag{23}
\]
By the definition of \(\overline I\), (23) is exactly \(\overline\Delta\in\overline I\).  Hence
\[
\Pr\{\overline\Delta\in\overline I\}\geq1-\alpha.
\tag{24}
\]
Finally, the same definition gives the deterministic identity
\[
\operatorname{diam}(\overline I)
=L^{-1}\sum_{\ell=1}^L\operatorname{diam}(I_\ell).
\tag{25}
\]
Taking expectations in (25), applying (21) to each of the fixed \(L=15\) pairs, and setting \(C_\alpha=C_{\alpha/L}\) proves
\[
\mathbb E\{\operatorname{diam}(\overline I)\}
\leq C_\alpha\left\{
\frac{\max(1,\sqrt{n\rho})}{n}+\eta
\right\}.
\]
\end{proof}
\par\smallskip\noindent \textit{Justification.} The fixed fifteen-pair target needs simultaneous rather than pointwise validity before interval averaging. \textit{Gap.} The published block bootstrap aggregates interior influence functions and does not guarantee finite-sample coverage at zero counts. \textit{Consumer.} The Buenos Aires-style fifteen-comparison DTT aggregate receives one causal-scale interval with a transparent multiplicity correction.

\begin{proposition}[prop:interior-reduction]\label{prop:interior-reduction}
Under the fixed interior probabilities \(r_{\ell j,n}=p^{\circ}_{\ell j}\in[\varepsilon_{\rm int},1-\varepsilon_{\rm int}]\), mutually independent Bernoulli role observations, and \(n_{\ell j}/n\to\pi_{\ell j}>0\), the outcome-independent sixty-role split satisfies, jointly over \(1\le\ell\le L\),
\[
\sqrt n(\widehat\Delta^E_{\ell}-\Delta_{\ell})
=\frac1{\sqrt n}\sum_{j=1}^{4}\sum_{a=1}^{n_{\ell j}}
\iota^{\circ}_{\ell j}(E_{\ell j a,n})+o_{P}(1)
\;\Rightarrow\;N(0,\Omega^{\circ}_{\ell}),
\]
with independent coordinates, where \(\iota^{\circ}_{\ell j}(e)=\gamma^{\circ}_{\ell j}(e-p^{\circ}_{\ell j})/\pi_{\ell j}\), \(\gamma^{\circ}_{\ell1}=\phi(v^{\circ}_{\ell})/\phi(q(p^{\circ}_{\ell1}))\), \(\gamma^{\circ}_{\ell2}=\phi(v^{\circ}_{\ell})/\phi(q(p^{\circ}_{\ell2}))\), \(\gamma^{\circ}_{\ell3}=-\phi(v^{\circ}_{\ell})/\phi(q(p^{\circ}_{\ell3}))\), \(\gamma^{\circ}_{\ell4}=-1\), \(v^{\circ}_{\ell}=q(p^{\circ}_{\ell1})+q(p^{\circ}_{\ell2})-q(p^{\circ}_{\ell3})\), and \(\Omega^{\circ}_{\ell}=\sum_{j=1}^{4}(\gamma^{\circ}_{\ell j})^2p^{\circ}_{\ell j}(1-p^{\circ}_{\ell j})/\pi_{\ell j}\). Consequently,
\[
\sqrt n\left\{L^{-1}\sum_{\ell=1}^{L}\widehat\Delta^E_{\ell}-\overline\Delta\right\}
\Rightarrow N\!\left(0,L^{-2}\sum_{\ell=1}^{L}\Omega^{\circ}_{\ell}\right).
\]
The interior margin is the finite-threshold standard-normal analogue of the regularity imposed in Djuazon and Tsyawo's interior analysis, but the displayed limit is proved independently for the mutually independent split-role experiment. It does not assert that this split experiment, or any rare triangular array, belongs to the sampling class of their Theorem 2.
\end{proposition}
\begin{proof}
Fix \(\ell\in\{1,\ldots,L\}\), and abbreviate \(p_j=p^{\circ}_{\ell j}\), \(n_j=n_{\ell j}\), and \(\pi_j=\pi_{\ell j}\). By \(\mathrm{ass:interior\mbox{-}probability\mbox{-}margin}\),
\[
 p=(p_1,p_2,p_3,p_4)\in[\varepsilon_{\rm int},1-\varepsilon_{\rm int}]^4\subset(0,1)^4.
\tag{1}
\]
Define the observed-law functional
\[
 g(a_1,a_2,a_3,a_4)=H(a_1,a_2,a_3)-a_4
 =\Phi\!\left(q(a_1)+q(a_2)-q(a_3)\right)-a_4,
 \qquad (a_1,a_2,a_3,a_4)\in(0,1)^4.
\tag{2}
\]
This definition does not define the causal effect. Instead, \(\mathrm{ass:interior\mbox{-}fixed\mbox{-}probabilities}\) gives \(r_{\ell j,n}=p_j\), and the already established \(\mathrm{prop:published\mbox{-}identification\mbox{-}baseline}\) therefore gives the separate identification equality
\[
 \Delta_\ell=g(p).
\tag{3}
\]

Because \(\Phi'(v)=\phi(v)>0\) for every finite \(v\), differentiating \(\Phi(q(a))=a\) yields
\[
 q'(a)=\frac{1}{\phi(q(a))},\qquad 0<a<1.
\tag{4}
\]
Thus \(q\), \(H\), and \(g\) are twice continuously differentiable on their respective open domains. Put
\[
 v=q(p_1)+q(p_2)-q(p_3).
\]
The chain rule and (4) imply
\[
 \nabla g(p)=\left(
 \frac{\phi(v)}{\phi(q(p_1))},
 \frac{\phi(v)}{\phi(q(p_2))},
 -\frac{\phi(v)}{\phi(q(p_3))},
 -1\right)=\gamma^{\circ}_{\ell}.
\tag{5}
\]
All divisions in (4)--(5) are valid because \(p_j\in(0,1)\) and \(\phi\) is strictly positive. By (1), there is a closed neighborhood \(U_\ell\subset(0,1)^4\) of \(p\) on which the Hessian of \(g\) is bounded.

By definition, \(K_{\ell j,n}=\sum_{a=1}^{n_j}E_{\ell j a,n}\) and
\[
 \widetilde r_{\ell j}=\frac{K_{\ell j,n}+1/2}{n_j+1}.
\]
Since \(0\le K_{\ell j,n}\le n_j\), this estimator is in \((0,1)\), including on the all-zero and all-one boundary samples, so every occurrence of \(g(\widetilde r_\ell)\) is well defined. Its centered decomposition is the exact identity
\[
 \widetilde r_{\ell j}-p_j
 =\frac{K_{\ell j,n}-n_jp_j}{n_j+1}
  +\frac{1/2-p_j}{n_j+1}.
\tag{6}
\]
The limit \(n_j/n\to\pi_j>0\) in \(\mathrm{ass:role\mbox{-}share\mbox{-}limits}\) implies \(n_j\asymp n\) and \(n_j\to\infty\). Under \(\mathrm{ass:bernoulli\mbox{-}role\mbox{-}sampling}\), the variables \(E_{\ell j a,n}\) have support \(\{0,1\}\), mean \(p_j\), and variance \(p_j(1-p_j)\). Hence
\[
 \operatorname{Var}(K_{\ell j,n}-n_jp_j)=n_jp_j(1-p_j),
 \qquad K_{\ell j,n}-n_jp_j=O_P(\sqrt{n_j})=O_P(\sqrt n),
\tag{7}
\]
where the stochastic order follows directly from Chebyshev's inequality. Equations (6)--(7), positivity of \(\pi_j\), and finiteness of the four roles yield
\[
 \widetilde r_\ell-p=O_P(n^{-1/2}).
\tag{8}
\]
Consequently, the event that the line segment between \(p\) and \(\widetilde r_\ell\) lies in \(U_\ell\) has probability tending to one. Taylor's formula with the bounded Hessian on that event, together with (3), (5), and (8), gives
\[
 \begin{aligned}
 \sqrt n\{\widehat\Delta^E_\ell-\Delta_\ell\}
 &=\sqrt n\{g(\widetilde r_\ell)-g(p)\}\\
 &=\sum_{j=1}^4\gamma^{\circ}_{\ell j}
   \frac{\sqrt n}{n_j+1}
   \sum_{a=1}^{n_j}(E_{\ell j a,n}-p_j)+o_P(1).
 \end{aligned}
\tag{9}
\]
Indeed, the deterministic add-half term in (6), after multiplication by \(\sqrt n\), is \(O(n^{-1/2})\), and the Taylor remainder is
\(
O_P(\sqrt n\,\|\widetilde r_\ell-p\|^2)=O_P(n^{-1/2})=o_P(1)
\).

The same positive-share limit gives
\[
 \frac{\sqrt n}{n_j+1}
 =\frac{1}{\sqrt n}\frac{n}{n_j+1}
 =\frac{1}{\sqrt n\,\pi_j}+o(n^{-1/2}).
\tag{10}
\]
Multiplying the last remainder by the centered count in (7) gives \(o_P(1)\). Substitution into (9), followed by the definition of \(\iota^{\circ}_{\ell j}\), proves the asserted asymptotic linear representation
\[
 \sqrt n\{\widehat\Delta^E_\ell-\Delta_\ell\}
 =\frac{1}{\sqrt n}\sum_{j=1}^4\sum_{a=1}^{n_j}
 \frac{\gamma^{\circ}_{\ell j}}{\pi_j}(E_{\ell j a,n}-p_j)+o_P(1)
 =\frac{1}{\sqrt n}\sum_{j=1}^4\sum_{a=1}^{n_j}
 \iota^{\circ}_{\ell j}(E_{\ell j a,n})+o_P(1).
\tag{11}
\]

Let \(S_{\ell n}\) denote the leading term in (11). Mutual independence from \(\mathrm{ass:disjoint\mbox{-}role\mbox{-}independence}\) and Bernoulli variances give
\[
 \operatorname{Var}(S_{\ell n})
 =\sum_{j=1}^4\frac{n_j}{n}
 \frac{(\gamma^{\circ}_{\ell j})^2p_j(1-p_j)}{\pi_j^2}
 \longrightarrow
 \sum_{j=1}^4\frac{(\gamma^{\circ}_{\ell j})^2p_j(1-p_j)}{\pi_j}
 =\Omega^{\circ}_\ell.
\tag{12}
\]
The limit is strictly positive: its \(j=4\) summand is \(p_4(1-p_4)/\pi_4>0\), because \(\gamma^{\circ}_{\ell4}=-1\), (1) holds, and \(\pi_4>0\). To obtain the limit law without importing any sampling-class claim, define
\[
 U_{\ell j,n}=\frac{K_{\ell j,n}-n_jp_j}{\sqrt{n_jp_j(1-p_j)}}.
\tag{13}
\]
For each fixed \((\ell,j)\), mutual independence and Bernoulli sampling imply \(K_{\ell j,n}\sim\operatorname{Binomial}(n_j,p_j)\); moreover, \(n_j\to\infty\) and \(p_j\in(0,1)\). The cited binomial central limit theorem \(\mathrm{lem:fixed\mbox{-}probability\mbox{-}binomial\mbox{-}clt}\) therefore gives \(U_{\ell j,n}\Rightarrow N(0,1)\). The entire collection in (13) is mutually independent for every \(n\). Hence, for any \(x=(x_{\ell j})\in\mathbb R^{4L}\), its joint distribution function factors and satisfies
\[
 \Pr\!\left(U_{\ell j,n}\le x_{\ell j}\text{ for every }\ell,j\right)
 =\prod_{\ell=1}^L\prod_{j=1}^4\Pr(U_{\ell j,n}\le x_{\ell j})
 \longrightarrow\prod_{\ell=1}^L\prod_{j=1}^4\Phi(x_{\ell j}).
\tag{14}
\]
The limiting distribution function is continuous, so (14) is joint convergence to \(4L\) independent standard-normal variables.

The leading term has the exact representation
\[
 S_{\ell n}=\sum_{j=1}^4
 \frac{\gamma^{\circ}_{\ell j}}{\pi_j}
 \sqrt{\frac{n_j}{n}p_j(1-p_j)}\,U_{\ell j,n}.
\tag{15}
\]
Every coefficient in (15) converges to
\(
\gamma^{\circ}_{\ell j}\sqrt{p_j(1-p_j)/\pi_j}
\).
Applying the corresponding deterministic linear map to (14) gives, jointly over the fixed \(L=15\) pairs,
\[
 (S_{1n},\ldots,S_{Ln})
 \Rightarrow N\!\left(0,\operatorname{diag}(\Omega^{\circ}_1,\ldots,\Omega^{\circ}_L)\right).
\tag{16}
\]
Thus the limiting coordinates are independent, and (16) also recovers the variance calculation (12). Because \(L\) is fixed, the vector of the remainders in (11) is \(o_P(1)\). It follows from (11) and (16) that the same joint limit holds for \(\bigl(\sqrt n(\widehat\Delta^E_\ell-\Delta_\ell)\bigr)_{\ell=1}^L\), proving the claimed pairwise limits with independent coordinates. Finally, the definition \(\overline\Delta=L^{-1}\sum_{\ell=1}^L\Delta_\ell\) and the deterministic linear averaging map give
\[
 \sqrt n\left\{L^{-1}\sum_{\ell=1}^{L}\widehat\Delta^E_{\ell}-\overline\Delta\right\}
 \Rightarrow N\!\left(0,L^{-2}\sum_{\ell=1}^{L}\Omega^{\circ}_{\ell}\right).
\]
Thus the causal interpretation comes solely from (3), while the expansion and both Gaussian limits are derived for the declared mutually independent Bernoulli split-role experiment under the fixed interior probabilities and positive role-share limits.
\end{proof}
\par\smallskip\noindent \textit{Justification.} The explicit normalization and influence function are proved directly for the mutually independent role split while preserving the published population DTT functional. \textit{Gap.} This is an independent fixed-threshold split-role analogue of interior Gaussian inference, not a claim that the split experiment or rare array lies in Djuazon and Tsyawo's Theorem 2 sampling class. \textit{Consumer.} The result supplies the regular fixed-probability benchmark against which the rare split-role limits are compared.

\begin{theorem}[thm:uniform-common-cutoff]\label{thm:uniform-common-cutoff}
There are universal finite constants \(M>0\) and \(\rho_0\in(0,1/16]\) such that, simultaneously for every \(1\le\ell\le L\), every \(0<\rho\le\rho_0\), and every \((c_{\ell1},c_{\ell2},c_{\ell3})\in[1/2,4]^3\),
\[
\left|u^2 C_{\rho,\ell}+1\right|\le \frac{M}{u^2}.
\]
The same bound holds on both exact cancellation planes \(x_\ell=0\) and \(z_\ell=0\).  Equivalently, uniformly on the complete coefficient box, \(C_{\rho,\ell}=-u^{-2}+O(u^{-4})\) as \(\rho\downarrow0\), without any division by \(A_\ell\).
\end{theorem}
\begin{proof}
Put \(R(v)=\Phi(-v)/\phi(v)\) for \(v>0\).  The change of variables \(w=v+s/v\), justified because the integrand is positive and integrable, gives the exact identity
\[
R(v)=\frac{1}{v}\int_0^\infty e^{-s}e^{-s^2/(2v^2)}\,ds.
\tag{1}
\]
For \(a\geq0\), Taylor's formula with Lagrange remainder gives
\[
e^{-a}=1-a+\frac{a^2}{2}+r(a),
\qquad |r(a)|\leq \frac{a^3}{6}.
\]
Using \(\int_0^\infty e^{-s}s^k\,ds=k!\) in (1), with \(a=s^2/(2v^2)\), yields
\[
R(v)=\frac{1}{v}\{1-v^{-2}+3v^{-4}+\varepsilon_v\},
\qquad |\varepsilon_v|\leq15v^{-6}.
\tag{2}
\]
For \(v\geq2\), the expression in braces is bounded away from zero.  Applying the exact algebraic identity
\[
(1+h)^{-1}=1-h+h^2-\frac{h^3}{1+h}
\]
to \(h=-v^{-2}+3v^{-4}+\varepsilon_v\) in (2) therefore gives an absolute finite constant \(K_0\) such that
\[
\left|m(-v)-\{v+v^{-1}-2v^{-3}\}\right|\leq K_0v^{-5}.
\tag{3}
\]
Differentiating the defining quotient \(m(x)=\phi(x)/\Phi(x)\) gives the exact identity \(m'(x)=-xm(x)-m(x)^2\).  Thus \(m'(-v)=-m(-v)\{m(-v)-v\}\), and substitution of (3), without differentiating its remainder, gives another absolute finite constant \(K_1\) such that
\[
\left|m'(-v)-\{-1+v^{-2}\}\right|\leq K_1v^{-4},
\qquad v\geq2.
\tag{4}
\]

We next obtain one cutoff for the entire coefficient box.  Let \(u=-q(\rho)\), let
\[
I=[-\log2,\log4],
\qquad v(a)=-q(\rho e^a),\quad a\in I,
\]
and choose \(u_0\geq4\) so large that \(\exp(3u_0^2/8)>4\).  Since (3), or directly (1), implies \(m(-w)>w\) for \(w>0\),
\[
\log\frac{\Phi(-u/2)}{\Phi(-u)}
=\int_{u/2}^{u}m(-w)\,dw
>\frac{3u^2}{8}.
\tag{5}
\]
Consequently, if \(u\geq u_0\), every probability in \([\rho/2,4\rho]\) has a negative normal quantile whose magnitude is at least \(u/2\).  The restriction \(\rho\leq1/16\) ensures that these probabilities lie in \((0,1)\).  The inverse-function rule is legitimate because \(\phi\) is strictly positive, and it gives
\[
v'(a)=-\frac{\rho e^a}{\phi(q(\rho e^a))}
=-\frac{1}{m(-v(a))}.
\]
Combining this formula with (5) and integrating along the interval between \(0\) and \(a\) gives the uniform estimate
\[
|v(a)-u|\leq\frac{2\log4}{u},
\qquad a\in I.
\tag{6}
\]

Fix a pair \(\ell\), suppress its index, and write \(a_j=\log c_j\) for \(j=1,2,3\).  For \(s,t\in[0,1]\), define
\[
a_s=(1-s)a_3+sa_1,
\qquad a_t=(1-t)a_3+ta_2.
\]
The coefficient-box assumption implies \(a_s,a_t,a_3\in I\).  Hence (6) shows, uniformly in \(s,t,\ell\), and the coefficients, that
\[
v(a_s)=u+O(u^{-1}),\qquad
v(a_t)=u+O(u^{-1}),\qquad
w_{s,t}:=v(a_s)+v(a_t)-v(a_3)=u+O(u^{-1}).
\tag{7}
\]
After increasing \(u_0\) if necessary, all three magnitudes in (7) are at least two.  Equations (3)--(4) then imply
\[
m(-v(a_s))=u+O(u^{-1}),\qquad
m(-v(a_t))=u+O(u^{-1}),
\]
and
\[
m'(-w_{s,t})=-1+O(u^{-2}).
\]
All constants are common over the stated box.  Since
\[
q(\rho e^{a_s})+q(\rho e^{a_t})-q(\rho e^{a_3})=-w_{s,t},
\]
division is valid because \(m\) is strictly positive, and elementary multiplication of the preceding uniform expansions gives a finite \(K_2\) such that
\[
\left|
\frac{m'(-w_{s,t})}{m(-v(a_s))m(-v(a_t))}+u^{-2}
\right|\leq K_2u^{-4}.
\tag{8}
\]
The definition of the mixed coefficient is exactly the integral of the ratio in (8) over the unit square.  Therefore
\[
B_{\rho,\ell}=-u^{-2}+R_{\rho,\ell},
\qquad |R_{\rho,\ell}|\leq K_2u^{-4},
\tag{9}
\]
uniformly over all \(L=15\) pairs.  Notice that no step divided by \(x_\ell\), \(z_\ell\), or \(A_\ell=x_\ell z_\ell\).

The same coefficient box gives \(|A_\ell|\leq(\log8)^2\).  For every real \(w\), the defining integral for \(\psi\) and the mean-value bound for the exponential give
\[
|\psi(w)-1|
\leq\int_0^1t|w|e^{t|w|}\,dt
\leq\frac12|w|e^{|w|}.
\tag{10}
\]
Equations (9)--(10) imply, uniformly,
\[
\psi(A_\ell B_{\rho,\ell})=1+O(u^{-2}),
\qquad
C_{\rho,\ell}=B_{\rho,\ell}\psi(A_\ell B_{\rho,\ell})
=-u^{-2}+O(u^{-4}).
\tag{11}
\]
Choose \(\rho_0=\min\{1/16,\Phi(-u_0)\}\), and choose \(M\) larger than the common absolute constant in (11).  Then
\[
|u^2C_{\rho,\ell}+1|\leq Mu^{-2}
\]
for every \(0<\rho\leq\rho_0\), every admissible coefficient triple, and every pair.  The rare-limit assumption guarantees \(u_n\to\infty\) along the triangular array.  Finally, if \(x_\ell=0\) or \(z_\ell=0\), then \(A_\ell=0\) and \(\psi(A_\ell B_{\rho,\ell})=1\); equations (8)--(11) are unchanged.  Thus the same certificate includes both exact cancellation planes.
\end{proof}

\begin{theorem}[thm:global-exact-factorization]\label{thm:global-exact-factorization}
For every \((p_1,p_2,p_3)\in(0,1)^3\), set \(b=p_3\), \(x=\log(p_1/p_3)\), \(z=\log(p_2/p_3)\), and \(A=xz\). Then
\[
\log H(p_1,p_2,p_3)-\log\!\left(\frac{p_1p_2}{p_3}\right)
=A B(p_1,p_2,p_3),
\qquad B(p_1,p_2,p_3)<0,
\]
and
\[
H(p_1,p_2,p_3)-\frac{p_1p_2}{p_3}
=\frac{p_1p_2}{p_3}A C(p_1,p_2,p_3),
\qquad C(p_1,p_2,p_3)<0.
\]
Consequently, for every \(p_4\in(0,1)\), the maps \(H(p_1,p_2,p_3)-p_4\) and \(P^{\rm P}(p_1,p_2,p_3)-p_4\) agree if and only if \(p_1=p_3\) or \(p_2=p_3\).
\end{theorem}
\begin{proof}
Fix \((p_1,p_2,p_3)\in(0,1)^3\), and define
\[
b=p_3,\qquad x=\log(p_1/p_3),\qquad z=\log(p_2/p_3),\qquad A=xz.
\]
Then \(be^x=p_1\) and \(be^z=p_2\).  For every \(s,t\in[0,1]\),
\[
be^{sx}=p_3^{1-s}p_1^s\in(0,1),
\qquad
be^{tz}=p_3^{1-t}p_2^t\in(0,1).
\tag{1}
\]
Thus all normal quantiles below are finite, while every inverse-Mills denominator is strictly positive by the definitions of \(m\), \(\phi\), and \(\Phi\).

For real \(a,d\) in the closed rectangle whose opposite corners are \((0,0)\) and \((x,z)\), define
\[
R_b(a,d)
=\log\Phi\!\left(q(be^a)+q(be^d)-q(b)\right)-\log b-a-d.
\tag{2}
\]
All arguments \(be^a\) and \(be^d\) lie in \((0,1)\) by the same geometric-interpolation calculation as in (1).  Since \(\Phi'=\phi>0\), the inverse-function rule and \(\Phi(q(p))=p\) give
\[
q'(p)=\frac{1}{\phi(q(p))},
\qquad
\frac{d}{da}q(be^a)
=\frac{be^a}{\phi(q(be^a))}
=\frac{1}{m(q(be^a))}.
\tag{3}
\]
Write \(S(a,d)=q(be^a)+q(be^d)-q(b)\).  The derivative of \(\log\Phi(v)\) is \(m(v)\).  Differentiating (2), first with respect to \(a\) and then with respect to \(d\), and applying (3) at each step yields
\[
\partial_{ad}R_b(a,d)
=\frac{m'(S(a,d))}
{m(q(be^a))m(q(be^d))}.
\tag{4}
\]
The definition of \(H\) and the identity \(\Phi(q(p))=p\) imply the exact boundary equalities
\[
R_b(a,0)=0,
\qquad
R_b(0,d)=0.
\tag{5}
\]
Indeed, the normal-link expression in (2) equals \(be^a\) when \(d=0\), and equals \(be^d\) when \(a=0\).

The integrand in (4) is continuous on the closed rectangle.  Applying the fundamental theorem of calculus twice to (4)--(5), using oriented integrals when \(x<0\) or \(z<0\), and then substituting \(a=sx\) and \(d=tz\), gives
\[
\begin{aligned}
R_b(x,z)
&=\int_0^x\!\int_0^z \partial_{ad}R_b(a,d)\,dd\,da\\
&=xz\int_0^1\!\int_0^1
\frac{m'\!\left(q(be^{sx})+q(be^{tz})-q(b)\right)}
{m(q(be^{sx}))m(q(be^{tz}))}\,dt\,ds\\
&=A B(p_1,p_2,p_3),
\end{aligned}
\tag{6}
\]
where the last equality is precisely the definition in \(\mathrm{def{:}mixed\text{-}coefficient}\).  This calculation also covers \(x=0\) or \(z=0\): in either case both sides of (6) are zero, and no step divides by \(x\), \(z\), or \(A\).

It remains to determine the sign.  Differentiating \(m(v)=\phi(v)/\Phi(v)\) and using \(\phi'(v)=-v\phi(v)\) gives
\[
m'(v)=-m(v)\{v+m(v)\}.
\tag{7}
\]
If \(v\geq0\), then \(v+m(v)>0\).  If \(v=-r<0\), where \(r>0\), symmetry of \(\phi\) and the substitution \(w=r+h\) give
\[
\Phi(-r)
=\int_r^\infty\phi(w)\,dw
=\phi(r)\int_0^\infty e^{-rh-h^2/2}\,dh
<\phi(r)\int_0^\infty e^{-rh}\,dh
=\frac{\phi(r)}{r}.
\tag{8}
\]
Consequently \(m(-r)=\phi(r)/\Phi(-r)>r\), so \(-r+m(-r)>0\).  Equation (7) therefore proves
\[
m'(v)<0\qquad\text{for every }v\in\mathbb R.
\tag{9}
\]
By (1), the integrand defining \(B(p_1,p_2,p_3)\) is continuous and strictly negative throughout the compact unit square, since both denominator factors are positive.  Hence
\[
B(p_1,p_2,p_3)<0.
\tag{10}
\]

By (2), \(R_b(x,z)=\log H(p_1,p_2,p_3)-\log(p_1p_2/p_3)\), because \(be^{x+z}=p_1p_2/p_3\).  Thus (6) proves the asserted logarithmic factorization.  Exponentiating it gives
\[
H(p_1,p_2,p_3)
=\frac{p_1p_2}{p_3}\exp\!\left(A B(p_1,p_2,p_3)\right).
\tag{11}
\]
For every \(w\in\mathbb R\), the definition \(\psi(w)=\int_0^1e^{tw}\,dt\) gives \(\psi(w)>0\), and the fundamental theorem of calculus gives
\[
e^w-1=w\psi(w),
\tag{12}
\]
including at \(w=0\).  Subtracting \(p_1p_2/p_3\) in (11), applying (12) with \(w=A B(p_1,p_2,p_3)\), and using the definition of \(C\) in \(\mathrm{def{:}mixed\text{-}coefficient}\), we obtain
\[
\begin{aligned}
H(p_1,p_2,p_3)-\frac{p_1p_2}{p_3}
&=\frac{p_1p_2}{p_3}A B(p_1,p_2,p_3)
\psi\!\left(A B(p_1,p_2,p_3)\right)\\
&=\frac{p_1p_2}{p_3}A C(p_1,p_2,p_3).
\end{aligned}
\tag{13}
\]
Equations (10) and \(\psi>0\) show that \(C(p_1,p_2,p_3)<0\).  Since \(p_1p_2/p_3>0\) and \(C(p_1,p_2,p_3)\neq0\), the right-hand side of (13) vanishes exactly when
\[
A=xz=0
\quad\Longleftrightarrow\quad
x=0\ \text{or}\ z=0
\quad\Longleftrightarrow\quad
p_1=p_3\ \text{or}\ p_2=p_3.
\tag{14}
\]
Finally, for any \(p_4\in(0,1)\), subtracting the same \(p_4\) from \(H(p_1,p_2,p_3)\) and \(P^{\rm P}(p_1,p_2,p_3)=p_1p_2/p_3\) does not change their difference.  The equivalence in (14) therefore proves the final assertion.
\end{proof}
\par\smallskip\noindent \textit{Justification.} The signed-rectangle identity exposes the full pointwise cancellation geometry without rarity or division by a cancellation coordinate. \textit{Gap.} Existing one-dimensional Mills results do not provide this mixed normal-link versus multiplicative identity or its exact two-plane zero set. \textit{Consumer.} The rare-panel theorem and its common-cutoff refinement evaluate this global coefficient along comparable-rarity sequences.

\begin{theorem}[thm:explicit-uniform-common-cutoff]\label{thm:explicit-uniform-common-cutoff}
The answer to \(\mathrm{oeq{:}uniform\text{-}common\text{-}cutoff}\) is affirmative.  One may take
\[
\rho_0=\min\!\left\{\frac1{16},\frac{\Phi(-6)}4\right\}
\qquad\text{and}\qquad M=221.
\]
For these constants, simultaneously for every \(1\le\ell\le L\), every \(0<\rho\le\rho_0\), and every \((c_{\ell1},c_{\ell2},c_{\ell3})\in[1/2,4]^3\),
\[
\left|u^2C_{\rho,\ell}+1\right|\le\frac{221}{u^2}.
\]
The inequality includes the planes \(x_\ell=0\) and \(z_\ell=0\).  Hence, uniformly on the whole coefficient box and along every sequence with \(\rho\downarrow0\),
\[
C_{\rho,\ell}=-u^{-2}+O(u^{-4}),
\]
including exact and drifting cancellation, and this conclusion requires no division by \(A_\ell\).
\end{theorem}
\begin{proof}
The exact factorization theorem identifies \(C_{\rho,\ell}\) as the complete-error coefficient and establishes that it is well defined on the stated probability domain.  It remains to control its size with one cutoff.  Fix \(0<\rho\le\rho_0\), suppress \(\ell\), and write \(c_j=c_{\ell j}\).  For \(a\in[1/2,4]\), define
\[
v_a=-q(\rho a).
\]
Since \(\rho a\le4\rho_0\le\Phi(-6)\) and \(q\) is strictly increasing, \(v_a\ge6\).  Also \(u=-q(\rho)>6\), because \(\rho\le\Phi(-6)/4<\Phi(-6)\).

We first compare every inverse tail with the same cutoff \(u\).  For \(t\) between \(0\) and \(\log a\), put \(v(t)=-q(\rho e^t)\).  The intermediate coefficient \(e^t\) lies in \([1/2,4]\), so \(v(t)\ge6\).  Differentiating the identity \(\Phi(q(p))=p\), and then using \(p=\rho e^t\), gives
\[
v'(t)=-\frac{\rho e^t}{\phi(q(\rho e^t))}
=-\frac1{m(-v(t))}.
\tag{5}
\]
The preceding lemma gives \(m(-v(t))>v(t)>0\).  Therefore
\[
\left|\frac{d}{dt}v(t)^2\right|
=\frac{2v(t)}{m(-v(t))}<2.
\]
Integration between \(0\) and \(\log a\) yields
\[
|v_a^2-u^2|\le2|\log a|,
\qquad
|v_a-u|
=\frac{|v_a^2-u^2|}{v_a+u}
\le\frac{2\log4}{u}<\frac4u.
\tag{6}
\]

Let \(x=\log(c_1/c_3)\), \(z=\log(c_2/c_3)\), and, for \(s,t\in[0,1]\), let
\[
a_s=c_3^{\,1-s}c_1^s,
\qquad
d_t=c_3^{\,1-t}c_2^t,
\qquad
v_s=v_{a_s},\quad \widetilde v_t=v_{d_t},\quad v_3=v_{c_3}.
\]
Geometric interpolation preserves \([1/2,4]\), so (6) applies to all three quantities, uniformly in \(s,t\).  In particular, if
\[
w=v_s+\widetilde v_t-v_3,
\]
then
\[
|w-u|<\frac{12}{u},
\qquad
w>u-\frac{12}{u}>\frac{2u}{3}>4.
\tag{7}
\]
The argument of \(m'\) in the integrand defining \(B_{\rho,\ell}\) is exactly
\[
q(\rho a_s)+q(\rho d_t)-q(\rho c_3)=-w.
\tag{8}
\]

Set \(D_s=m(-v_s)\) and \(\widetilde D_t=m(-\widetilde v_t)\).  From (6), \(v_s>u-4/u>8u/9\), and the same inequality holds for \(\widetilde v_t\).  The lemma and (6) now imply
\[
|D_s-u|
\le |D_s-v_s|+|v_s-u|
<\frac2{v_s}+\frac4u
<\frac7u,
\tag{9}
\]
with the identical bound for \(\widetilde D_t\).  Write
\[
D_s=u(1+e_s),\qquad \widetilde D_t=u(1+\widetilde e_t).
\]
Then (9) and \(u>6\) give
\[
|e_s|,|\widetilde e_t|<d:=\frac7{u^2}<\frac15.
\tag{10}
\]
For any \(|e|,|f|\le d<1/5\),
\[
\left|\frac1{(1+e)(1+f)}-1\right|
\le\frac{2d+d^2}{(1-d)^2}<4d,
\qquad
\frac1{|(1+e)(1+f)|}<2.
\tag{11}
\]
By (7) and the lemma, there is an \(e_0=e_0(s,t)\) such that
\[
m'(-w)=-1+e_0,
\qquad
|e_0|\le\frac7{w^2}<\frac{16}{u^2}.
\tag{12}
\]
Combining (8)--(12), the mixed-derivative integrand \(K(s,t)\) in the definition of \(B_{\rho,\ell}\) satisfies
\[
\begin{aligned}
\left|K(s,t)+\frac1{u^2}\right|
&=\frac1{u^2}
\left|1-\frac1{(1+e_s)(1+\widetilde e_t)}
+\frac{e_0}{(1+e_s)(1+\widetilde e_t)}\right|\\
&\le\frac1{u^2}\left(\frac{28}{u^2}+\frac{32}{u^2}\right)
=\frac{60}{u^4}.
\end{aligned}
\tag{13}
\]
All constants in (13) are independent of \(s,t,\ell\), and the coefficients.  Integrating over \([0,1]^2\) gives
\[
\left|B_{\rho,\ell}+\frac1{u^2}\right|\le\frac{60}{u^4},
\qquad
|B_{\rho,\ell}|\le\frac3{u^2}.
\tag{14}
\]

Because \(c_1/c_3,c_2/c_3\in[1/8,8]\), one has \(|x|,|z|\le\log8<3\), whence \(|A_\ell|=|xz|<9\).  Thus (14) implies
\[
|A_\ell B_{\rho,\ell}|\le\frac{27}{u^2}<\frac34.
\tag{15}
\]
For every real \(y\) with \(|y|<3/4\), the integral definition of \(\psi\) and the mean-value theorem give
\[
|\psi(y)|\le e^{|y|}<3,
\qquad
|\psi(y)-1|
\le\int_0^1 t|y|e^{t|y|}\,dt
\le\frac32|y|.
\tag{16}
\]
Apply (16) with \(y=A_\ell B_{\rho,\ell}\).  Equations (15)--(16) yield
\[
\left|\psi(A_\ell B_{\rho,\ell})-1\right|
\le\frac{81}{2u^2}<\frac{41}{u^2}.
\tag{17}
\]
Using \(C_{\rho,\ell}=B_{\rho,\ell}\psi(A_\ell B_{\rho,\ell})\), (14), (16), and (17), we conclude that
\[
\begin{aligned}
|u^2C_{\rho,\ell}+1|
&\le |u^2B_{\rho,\ell}+1|\,
|\psi(A_\ell B_{\rho,\ell})|
+|1-\psi(A_\ell B_{\rho,\ell})|\\
&\le\frac{180}{u^2}+\frac{41}{u^2}
=\frac{221}{u^2}.
\end{aligned}
\tag{18}
\]
No step divided by \(A_\ell\).  If \(x_\ell=0\) or \(z_\ell=0\), then \(A_\ell=0\) and \(\psi(A_\ell B_{\rho,\ell})=\psi(0)=1\); equations (5)--(18) remain valid.  Thus (18) covers both exact cancellation planes.  Finally, the assumed limit \(\rho_n\to0\) places every admissible sequence below \(\rho_0\) eventually, so (18) is equivalently the asserted uniform expansion along the rare limit.
\end{proof}

\begin{lemma}[lem:quantitative-left-mills-bounds]\label{lem:quantitative-left-mills-bounds}
For every \(v\ge 2\), define
\[
I(v)=\int_0^\infty e^{-s}e^{-s^2/(2v^2)}\,ds.
\]
Then
\[
\frac{\Phi(-v)}{\phi(v)}=\frac{I(v)}{v},\qquad
0<m(-v)-v\le \frac{2}{v},
\qquad
\left|m'(-v)+1\right|\le \frac{7}{v^2}.
\]
\end{lemma}
\begin{proof}
Substitute \(w=v+s/v\) in \(\Phi(-v)=\int_v^\infty\phi(w)\,dw\).  Since
\[
\frac{\phi(v+s/v)}{\phi(v)}
=\exp\!\left(-s-\frac{s^2}{2v^2}\right),
\]
the substitution gives the asserted identity.  Put \(\delta(v)=1-I(v)\) and \(a=s^2/(2v^2)\).  The elementary inequalities
\[
0\le 1-e^{-a}\le a,
\qquad
0\le a-(1-e^{-a})\le \frac{a^2}{2}
\tag{1}
\]
hold for every \(a\ge0\); they follow, respectively, by integrating \(e^{-t}\le1\) and by applying Taylor's theorem to \(e^{-a}\) through its quadratic term.  Repeated integration by parts gives
\[
\int_0^\infty e^{-s}s^2\,ds=2,
\qquad
\int_0^\infty e^{-s}s^4\,ds=24.
\]
Consequently, (1) yields
\[
0<\delta(v)\le \frac1{v^2},
\qquad
0\le 1-v^2\delta(v)\le \frac3{v^2}.
\tag{2}
\]
For \(v\ge2\), (2) implies \(I(v)\ge3/4\).  Because \(m(-v)=v/I(v)\),
\[
0<m(-v)-v=\frac{v\delta(v)}{I(v)}
\le \frac{4}{3v}<\frac2v.
\tag{3}
\]
Direct differentiation of \(m(w)=\phi(w)/\Phi(w)\) gives
\[
m'(w)=-m(w)\{w+m(w)\}.
\tag{4}
\]
Let \(X=v^2\delta(v)\).  Equations (2)--(4) show that \(0<X\le1\) and
\[
-m'(-v)=m(-v)\{m(-v)-v\}=\frac{X}{I(v)^2}.
\]
Furthermore,
\[
I(v)^{-2}-1
=\frac{\delta(v)\{1+I(v)\}}{I(v)^2}
\le \frac{32}{9v^2}.
\]
It follows that
\[
\left|-m'(-v)-1\right|
\le |X-1|+X\{I(v)^{-2}-1\}
\le \frac3{v^2}+\frac{32}{9v^2}
<\frac7{v^2},
\]
which proves all three bounds.
\end{proof}

\begin{lemma}[lem:fixed-probability-binomial-clt]\label{lem:fixed-probability-binomial-clt}
Let \(p\in(0,1)\) be fixed. If \(X_1,X_2,\ldots\) are independent \(\operatorname{Bernoulli}(p)\) random variables and \(K_m=\sum_{a=1}^mX_a\), then
\[
 \frac{K_m-mp}{\sqrt{mp(1-p)}}\Rightarrow N(0,1)
 \qquad\text{as }m\to\infty.
\]
\end{lemma}
\begin{proof}
Put \(q=1-p\) and \(\sigma=\sqrt{pq}\).  The hypotheses \(0<p<1\) imply \(q>0\) and \(\sigma>0\), so the standardized variables
\[
 Z_a=\frac{X_a-p}{\sigma}
\]
are well defined.  Independence of the \(X_a\)'s implies independence of the \(Z_a\)'s, and the Bernoulli law gives
\[
 \mathbb E Z_a=0,
 \qquad
 \mathbb E Z_a^2
 =\frac{p(1-p)^2+(1-p)p^2}{p(1-p)}=1.
\]
Moreover, \(Z_a\) takes only the two finite values \(q/\sigma\) and \(-p/\sigma\); hence
\[
 \mu_3:=\mathbb E|Z_a|^3<\infty.
\]

We first calculate the limiting characteristic function without invoking a central limit theorem.  Taylor's formula with integral remainder, applied to \(v\mapsto e^{ivy}\) on \([0,1]\), gives, for every real \(y\),
\[
 e^{iy}=1+iy-\frac{y^2}{2}+r(y),
 \qquad
 r(y)=\frac{(iy)^3}{2}\int_0^1(1-v)^2e^{ivy}\,dv.
\]
Since \(|e^{ivy}|=1\),
\[
 |r(y)|\leq \frac{|y|^3}{2}\int_0^1(1-v)^2\,dv
 =\frac{|y|^3}{6}.
\]
Therefore the characteristic function \(\chi(s)=\mathbb E e^{isZ_1}\) satisfies
\[
 \chi(s)=1-\frac{s^2}{2}+R(s),
 \qquad
 |R(s)|\leq \frac{\mu_3|s|^3}{6},
\]
where the linear term vanishes because \(\mathbb E Z_1=0\), and the quadratic term equals \(-s^2/2\) because \(\mathbb E Z_1^2=1\).

Define
\[
 S_m=\frac{1}{\sqrt m}\sum_{a=1}^m Z_a
 =\frac{K_m-mp}{\sqrt{mp(1-p)}}.
\]
By independence, its characteristic function is
\[
 \varphi_m(t)=\mathbb E e^{itS_m}
 =\left\{\chi\!\left(\frac{t}{\sqrt m}\right)\right\}^{m}.
\]
For fixed \(t\in\mathbb R\), write
\[
 a_m(t)=\chi\!\left(\frac{t}{\sqrt m}\right)-1
 =-\frac{t^2}{2m}+R\!\left(\frac{t}{\sqrt m}\right).
\]
The preceding remainder bound yields
\[
 \left|a_m(t)+\frac{t^2}{2m}\right|
 \leq \frac{\mu_3|t|^3}{6m^{3/2}},
\]
so \(a_m(t)=O_t(m^{-1})\) and \(|a_m(t)|<1/2\) for all sufficiently large \(m\).  On \(|z|<1\), the absolutely convergent power series for the logarithm gives
\[
 \left|\log(1+z)-z\right|
 \leq \sum_{k=2}^{\infty}\frac{|z|^k}{k}
 \leq \frac{|z|^2}{2(1-|z|)}.
\]
Consequently,
\[
 m\log\{1+a_m(t)\}
 =m a_m(t)+O_t\!\left(m|a_m(t)|^2\right)
 =-\frac{t^2}{2}+O_t(m^{-1/2}),
\]
and exponentiation shows
\[
 \varphi_m(t)\longrightarrow e^{-t^2/2}
 \qquad (t\in\mathbb R).
\]

For completeness, we now derive distributional convergence directly from this characteristic-function limit.  Write \(\phi(v)=(2\pi)^{-1/2}e^{-v^2/2}\).  Let \(G\) be a random variable with density \(\phi\), independent of every \(S_m\), and fix \(\varepsilon>0\).  Set \(U_{m,\varepsilon}=S_m+\varepsilon G\).  The elementary Gaussian Fourier identity
\[
 \frac{1}{\varepsilon}\phi\!\left(\frac{x-y}{\varepsilon}\right)
 =\frac{1}{2\pi}\int_{\mathbb R}
 e^{-it(x-y)}e^{-\varepsilon^2t^2/2}\,dt
\]
follows by completing the square in the Gaussian integral.  Conditioning on \(S_m\), and using Fubini's theorem with the integrable dominating function \(e^{-\varepsilon^2t^2/2}\), shows that \(U_{m,\varepsilon}\) has density
\[
 g_{m,\varepsilon}(x)=\frac{1}{2\pi}\int_{\mathbb R}
 e^{-itx}\varphi_m(t)e^{-\varepsilon^2t^2/2}\,dt.
\]
Because \(|\varphi_m(t)|\leq1\), dominated convergence and the already proved pointwise limit imply
\[
 \sup_{x\in\mathbb R}|g_{m,\varepsilon}(x)-g_{\varepsilon}(x)|
 \leq \frac{1}{2\pi}\int_{\mathbb R}
 |\varphi_m(t)-e^{-t^2/2}|e^{-\varepsilon^2t^2/2}\,dt
 \longrightarrow0,
\]
where another application of the Gaussian Fourier identity gives
\[
 g_{\varepsilon}(x)
 =\frac{1}{\sqrt{1+\varepsilon^2}}
 \phi\!\left(\frac{x}{\sqrt{1+\varepsilon^2}}\right).
\]
Also \(\mathbb E U_{m,\varepsilon}=0\) and, by independence,
\[
 \mathbb E U_{m,\varepsilon}^2
 =\mathbb E S_m^2+\varepsilon^2\mathbb E G^2
 =1+\varepsilon^2.
\]
The same second moment holds for the centered normal variable having density \(g_{\varepsilon}\).  Thus, for fixed \(x\) and any \(R>|x|\), uniform convergence of the densities on \([-R,x]\) and the bound \(\Pr(|U_{m,\varepsilon}|\geq R)\leq(1+\varepsilon^2)/R^2\), which follows from \(R^2\mathbf 1\{|U_{m,\varepsilon}|\geq R\}\leq U_{m,\varepsilon}^2\), imply
\[
 \Pr(U_{m,\varepsilon}\leq x)
 \longrightarrow
 \Phi\!\left(\frac{x}{\sqrt{1+\varepsilon^2}}\right).
\]
Indeed, after truncation at \(-R\), the integral of the density difference is at most \((R+|x|)\sup_y|g_{m,\varepsilon}(y)-g_{\varepsilon}(y)|\); the two omitted left tails are together at most \(2(1+\varepsilon^2)/R^2\).  First let \(m\to\infty\), and then let \(R\to\infty\).

It remains to remove the auxiliary smoothing.  For every \(\delta>0\), the inclusions obtained on the event \(\{|\varepsilon G|\leq\delta\}\) give
\[
 \Pr(U_{m,\varepsilon}\leq x-\delta)-\Pr(|\varepsilon G|>\delta)
 \leq \Pr(S_m\leq x)
 \leq \Pr(U_{m,\varepsilon}\leq x+\delta)+\Pr(|\varepsilon G|>\delta).
\]
For \(0<\varepsilon<1\), choose \(\delta=\sqrt\varepsilon\).  Since \(\mathbb E G^2=1\),
\[
 \Pr(|\varepsilon G|>\sqrt\varepsilon)
 =\Pr(|G|>\varepsilon^{-1/2})\leq\varepsilon.
\]
Taking first the lower and upper limits as \(m\to\infty\), and then letting \(\varepsilon\downarrow0\), continuity of \(\Phi\) yields
\[
 \liminf_{m\to\infty}\Pr(S_m\leq x)
 \geq\Phi(x),
 \qquad
 \limsup_{m\to\infty}\Pr(S_m\leq x)
 \leq\Phi(x).
\]
Hence \(\Pr(S_m\leq x)\to\Phi(x)\) for every real \(x\), which is precisely
\[
 \frac{K_m-mp}{\sqrt{mp(1-p)}}\Rightarrow N(0,1).
\]
\end{proof}

\begin{lemma}[lem:fixed-compact-gaussian-tail-transfer]\label{lem:fixed-compact-gaussian-tail-transfer}
Fix deterministic constants \(0<d\leq D<\infty\). There are constants \(\varepsilon _0\in(0,1/D)\) and \(K<\infty\), depending only on \(d,D\), such that the following assertions hold for every \(0<\varepsilon\leq\varepsilon _0\). Put \(v=-q(\varepsilon)\). Uniformly for \(a\in[d,D]\),
\[
 \left|q(\varepsilon a)+v-\frac{\log a}{v}\right|
 \leq \frac{K}{v^3},
 \qquad
 \left|\frac{m(q(\varepsilon a))}{v}-1\right|
 \leq \frac{K}{v^2}.
\]
For every \((a_1,a_2,a_3)\in[d,D]^3\), let \(Q(a)=a_1a_2/a_3\) and \(p_j=\varepsilon a_j\). Then
\[
 \left|\frac{H(p_1,p_2,p_3)}{\varepsilon Q(a)}-1\right|
 \leq\frac{K}{v^2},
\]
\[
 \left|\partial_1H(p)-\frac{a_2}{a_3}\right|
 +\left|\partial_2H(p)-\frac{a_1}{a_3}\right|
 +\left|\partial_3H(p)+\frac{a_1a_2}{a_3^2}\right|
 \leq\frac{K}{v^2},
 \qquad
 \max_{1\leq j,k\leq3}|\partial_{jk}H(p)|\leq\frac{K}{\varepsilon}.
\]
Thus the value, gradient, and Hessian transfers are uniform on every fixed compact coefficient interval, with one cutoff for the whole interval.
\end{lemma}
\begin{proof}
Write
\[
 R(w)=\frac{\Phi(-w)}{\phi(w)},\qquad w>0.
\]
Changing variables first from \(y\) to \(t=y-w\), and then from \(t\) to \(s=wt\), gives the exact identity
\[
 R(w)=\int_0^\infty e^{-wt-t^2/2}\,dt
 =\frac1w\int_0^\infty e^{-s}e^{-s^2/(2w^2)}\,ds.
\tag{1}
\]
For \(x\geq0\), Taylor's formula with integral remainder gives
\[
 0\leq e^{-x}-(1-x)\leq \frac{x^2}{2}.
\]
Since \(\int_0^\infty e^{-s}s^k\,ds=k!\) for nonnegative integers \(k\), equation (1) implies
\[
 R(w)=\frac1w\left(1-\frac1{w^2}+r_w\right),
 \qquad 0\leq r_w\leq\frac3{w^4}.
\tag{2}
\]
For all sufficiently large \(w\), the bracket in (2) is bounded below by \(1/2\). Algebraic inversion of (2) then yields a universal finite constant \(K_0\) such that
\[
 m(-w)=\frac1{R(w)}=w+\frac1w+e_w,
 \qquad |e_w|\leq\frac{K_0}{w^3}.
\tag{3}
\]
In particular, after increasing the deterministic cutoff if necessary,
\[
 w\leq m(-w)\leq2w.
\tag{4}
\]
Every division in (1)--(4) is valid because \(w>0\), \(\phi(w)>0\), and \(\Phi(-w)>0\).

Let \(v=-q(\varepsilon)\) and, for a positive coefficient \(a\), let \(w_a=-q(\varepsilon a)\). Since \(\Phi(-M)>0\) for every finite \(M\), the relation \(\Phi(-v)=\varepsilon\) implies \(v\to\infty\) as \(\varepsilon\downarrow0\). Uniformly for \(a\) in any fixed compact subset of \((0,\infty)\), \(\varepsilon a\to0\), so also \(w_a\to\infty\). Differentiating \(\log\Phi(-w)\), and using \(m=\phi/\Phi\), gives
\[
 \log a
 =\log\Phi(-w_a)-\log\Phi(-v)
 =\int_{w_a}^{v}m(-s)\,ds.
\tag{5}
\]
Let \(K_1=\max\{|\log d|,|\log D|\}\). For sufficiently small \(\varepsilon\), equation (4) and (5) rule out \(w_a<v/2\) and \(w_a>2v\), uniformly for \(a\in[d,D]\): either alternative would make the absolute value of the integral at least a positive constant times \(v^2\), while \(|\log a|\leq K_1\). Thus \(w_a\in[v/2,2v]\). Equation (4), applied on the interval with endpoints \(v,w_a\), now gives
\[
 |v-w_a|\leq\frac{2K_1}{v}.
\tag{6}
\]
Put \(\delta_a=v-w_a\). Substitution of (3) into (5), with oriented integration when \(\delta_a<0\), gives
\[
 \begin{aligned}
 \log a
 &=\int_{v-\delta_a}^{v}
 \left(s+\frac1s+O(s^{-3})\right)ds\\
 &=v\delta_a+O(v^{-2}),
 \end{aligned}
\tag{7}
\]
uniformly for \(a\in[d,D]\); here (6) bounds the quadratic term \(\delta_a^2/2\), the integral of \(s^{-1}\), and the integrated remainder in (3). Dividing (7) by \(v>0\) proves
\[
 q(\varepsilon a)=-v+\frac{\log a}{v}+O(v^{-3})
\tag{8}
\]
uniformly on \([d,D]\). Equations (3), (6), and (8) also give
\[
 m(q(\varepsilon a))=v+O(v^{-1})
\tag{9}
\]
uniformly on that interval. The constants hidden in (7)--(9) depend only on its two fixed endpoints, so they can be absorbed into a single \(K\).

Now fix \((a_1,a_2,a_3)\in[d,D]^3\), put \(q_j=q(\varepsilon a_j)\), and define
\[
 Q=\frac{a_1a_2}{a_3},
 \qquad s=q_1+q_2-q_3.
\]
The coefficient \(Q\) belongs to the deterministic compact interval
\[
 \left[\frac{d^2}{D},\frac{D^2}{d}\right]\subset(0,\infty).
\tag{10}
\]
Decrease \(\varepsilon_0\), if necessary, so that \(\varepsilon D^2/d<1\); then \(q(\varepsilon Q)\) is defined for every coefficient triple under consideration.
Apply (8) both on \([d,D]\) and on the fixed interval in (10). Since \(\log Q=\log a_1+\log a_2-\log a_3\), this gives
\[
 s=-v+\frac{\log Q}{v}+O(v^{-3}),
 \qquad
 q(\varepsilon Q)=-v+\frac{\log Q}{v}+O(v^{-3}),
\tag{11}
\]
uniformly on \([d,D]^3\). Hence \(|s-q(\varepsilon Q)|=O(v^{-3})\). Every point between these two arguments has the form \(-v+O(v^{-1})\), and (3) therefore bounds its inverse Mills ratio by a fixed multiple of \(v\). The fundamental theorem of calculus applied to \(\log\Phi\), whose derivative is \(m\), consequently yields
\[
 \left|\log\frac{\Phi(s)}{\varepsilon Q}\right|
 =\left|\int_{q(\varepsilon Q)}^s m(y)\,dy\right|
 =O(v^{-2}).
\tag{12}
\]
After making \(\varepsilon_0\) smaller so that the last bound has absolute value at most one, the inequality \(|e^t-1|\leq e|t|\) for \(|t|\leq1\) and the definition of \(H\) prove
\[
 \frac{H(\varepsilon a_1,\varepsilon a_2,\varepsilon a_3)}{\varepsilon Q}
 =1+O(v^{-2}).
\tag{13}
\]

It remains to control derivatives. Differentiating \(\Phi(q(p))=p\), using \(\Phi'=\phi>0\), gives \(q'(p)=1/\phi(q(p))\). With \((e_1,e_2,e_3)=(1,1,-1)\), direct differentiation therefore gives
\[
 \partial_jH(p)=e_j\frac{\phi(s)}{\phi(q_j)}
 =e_j\frac{H(p)m(s)}{\varepsilon a_jm(q_j)},
\tag{14}
\]
and
\[
 \partial_{jk}H(p)=\partial_jH(p)
 \left\{-\frac{s e_k}{\phi(q_k)}
 +\boldsymbol 1\{j=k\}\frac{q_j}{\phi(q_j)}\right\}.
\tag{15}
\]
Equations (3), (9), and (11) show uniformly that
\[
 m(s)=v+O(v^{-1}),
 \qquad \frac{m(s)}{m(q_j)}=1+O(v^{-2}).
\tag{16}
\]
Substitution of (13) and (16) into (14) gives
\[
 \partial_1H(p)=\frac{a_2}{a_3}+O(v^{-2}),\qquad
 \partial_2H(p)=\frac{a_1}{a_3}+O(v^{-2}),\qquad
 \partial_3H(p)=-\frac{a_1a_2}{a_3^2}+O(v^{-2}).
\tag{17}
\]
Finally, (9) and \(\phi(q_j)=\Phi(q_j)m(q_j)=\varepsilon a_jm(q_j)\) imply
\[
 \phi(q_j)\geq c_0\varepsilon v,
 \qquad |q_j|+|s|\leq C_0v
\tag{18}
\]
for finite positive constants \(c_0,C_0\) depending only on \(d,D\). Equation (17) bounds every first derivative, so (15)--(18) give
\[
 \max_{j,k}|\partial_{jk}H(p)|\leq\frac{K}{\varepsilon}.
\]
Taking one sufficiently small common \(\varepsilon_0\) and one sufficiently large common \(K\) proves every asserted inequality simultaneously.
\end{proof}

\section*{Technical internal limitation (diagnostic only)}

The literal statement that every all-full table with all \(n_j\le6\) must yield an empty retained list is false without a calibration restriction: for example, with \(n_j=1\) and \(\beta=1/2\), the admissible point \(p_j=1/4\) has deviance \(8\log4<16=8/\beta\). The formal theorem therefore states the exact sufficient inequality for empty-list behavior. This diagnostic does not affect all-zero retention or coverage. The global Lipschitz constant is deliberately enormous, so finite termination is a mathematical certificate rather than a practical runtime guarantee.

\section{Honest scope}

The pointwise cancellation theorem is a population result for every \(p_1,p_2,p_3\in(0,1)\); it needs neither rarity, sampling, an interior margin, nor the coefficient box. The uniform expansion \(C_{\rho,\ell}=-u^{-2}+O(u^{-4})\) is narrower: it uses the common-rarity evaluations on the fixed comparable-rarity box and the rare limit, while remaining valid on both exact cancellation planes. The Gaussian, product-Poisson, and rational confidence-set results retain the predeclared mutually independent split-role design and do not establish covariance formulas or joint rare-event limits for overlapping original blocks. Fixed interior probabilities are handled by an independently proved split-role analogue of the published interior result, not by asserting that the rare triangular array belongs to Djuazon and Tsyawo's Theorem 2 class. The conservative aggregate interval must include zero in the published 37-treated-block application, so no finite-sample power claim is made. The global Lipschitz certificate ensures finite termination but is intentionally conservative computationally.

\begin{thebibliography}{99}

  \bibitem{DjuazonTsyawo2026} Nelly K. Djuazon and Emmanuel Selorm Tsyawo (2026), ``Quantile and Distribution Treatment Effects on the Treated with Possibly Non-Continuous Outcomes,'' arXiv:2408.07842v2.

  \bibitem{Pinelis2019InverseMills} Iosif Pinelis (2019), ``Exact Bounds on the Inverse Mills Ratio and Its Derivatives,'' Complex Analysis and Operator Theory 13, 1643--1651; arXiv:1512.00120.

  \bibitem{Baricz2008Mills} Árpád Baricz (2008), ``Mills' Ratio: Monotonicity Patterns and Functional Inequalities,'' Journal of Mathematical Analysis and Applications 340, 1362--1370.

  \bibitem{ChenFang2019} Qihui Chen and Zheng Fang (2019), ``Inference on Functionals under First Order Degeneracy,'' Journal of Econometrics 210, 459--481; arXiv:1901.04861.

  \bibitem{Marohn2001} Frank Marohn (2001), ``Limits of Truncation Experiments,'' Probability and Mathematical Statistics 21, 71--88.

  \bibitem{SachsGabrielFay2026} Michael C. Sachs, Erin E. Gabriel, and Michael P. Fay (2026), ``Exact Confidence Intervals for Functions of Parameters in the \(k\)-Sample Multinomial Problem,'' arXiv:2406.19141.

  \bibitem{BrownCaiDasGupta2001} Lawrence D. Brown, T. Tony Cai, and Anirban DasGupta (2001), ``Interval Estimation for a Binomial Proportion,'' Statistical Science 16, 101--133.

  \bibitem{FernandezValEtAl2024} Iván Fernández-Val, Jonas Meier, Aico van Vuuren, and Francis Vella (2024), ``Distribution Regression Difference-in-Differences,'' arXiv:2409.02311.

  \bibitem{RothSantAnna2023} Jonathan Roth and Pedro H. C. Sant'Anna (2023), ``When Is Parallel Trends Sensitive to Functional Form?'' Econometrica 91, 737--747.

  \bibitem{Wooldridge2023} Jeffrey M. Wooldridge (2023), ``Simple Approaches to Nonlinear Difference-in-Differences with Panel Data,'' The Econometrics Journal 26, C31--C66.

  \bibitem{SchaferStark2009} Chad M. Schafer and Philip B. Stark (2009), ``Constructing Confidence Regions of Optimal Expected Size,'' Journal of the American Statistical Association 104, 1080--1089.

  \bibitem{ArratiaGoldsteinGordon1989} Richard Arratia, Larry Goldstein, and Louis Gordon (1989), ``Two Moments Suffice for Poisson Approximations: The Chen--Stein Method,'' Annals of Probability 17, 9--25.

\end{thebibliography}

\end{document}